\documentclass[11pt]{article}

\usepackage[T1]{fontenc}
\usepackage[utf8]{inputenc}
\usepackage{lmodern}
\usepackage[margin=1in]{geometry}
\usepackage{amsmath,amssymb,amsthm}
\usepackage{booktabs,longtable,tabularx,array}
\usepackage{microtype}
\usepackage[hidelinks]{hyperref}

\setlength{\parindent}{0pt}
\setlength{\parskip}{0.55em}
\setlength{\emergencystretch}{3em}
\urlstyle{tt}

\newcommand{\F}{\mathbb F}
\newcommand{\N}{\mathbb N}
\newcommand{\IR}{\mathsf{IR}}
\newcommand{\Time}{\operatorname{TIME}}
\newcommand{\Eval}{\operatorname{Eval}}
\newcommand{\Quote}{\operatorname{Quote}}
\newcommand{\Specialize}{\operatorname{Specialize}}
\newcommand{\Compress}{\mathsf{Compress}}
\newcommand{\Introspect}{\mathsf{Introspect}}
\newcommand{\AnswerReduce}{\mathsf{AnswerReduce}}
\newcommand{\Repeat}{\mathsf{Repeat}}
\newcommand{\ComputeIntroVerifier}{\mathsf{ComputeIntroVerifier}}
\newcommand{\ComputeAnsVerifier}{\mathsf{ComputeAnsVerifier}}
\newcommand{\ComputeParrepVerifier}{\mathsf{ComputeParrepVerifier}}
\newcommand{\ComputeSampler}{\mathsf{ComputeSampler}}
\newcommand{\poly}{\operatorname{poly}}
\newcommand{\code}[1]{\nolinkurl{#1}}
\newcommand{\gt}[2]{\footnote{Ground truth: \code{#1} in \code{#2}.}}
\newcommand{\grade}[1]{\textsc{\MakeLowercase{#1}}}
\newcommand{\constructed}{\grade{CONSTRUCTED}}
\newcommand{\claim}[2]{\code{#1}: \textsc{#2}}
\newcolumntype{P}[1]{>{\raggedright\arraybackslash}p{#1}}

\newtheorem{theorem}{Theorem}[section]
\newtheorem{lemma}[theorem]{Lemma}
\newtheorem{proposition}[theorem]{Proposition}
\theoremstyle{definition}
\newtheorem{definition}[theorem]{Definition}
\newtheorem{remark}[theorem]{Remark}

\newenvironment{reminder}{%
  \par\smallskip\noindent\begin{minipage}{\linewidth}%
  \hrule\medskip\small\textbf{Reminder. }\ignorespaces
}{%
  \par\medskip\hrule\end{minipage}\smallskip\par
}

\title{Analytic Underpinnings of the Symbolic Answer-Reduction Pipeline}
\author{}
\date{4 September 2026}

\begin{document}
\maketitle

\begin{abstract}
The objects manipulated by answer reduction are finite descriptions of
machines, circuits, formulas, polynomials, samplers, and verifiers.  The Julia
prototype realizes part of the corresponding analytic chain as typed
transformations of explicit intermediate representations.  It constructs
finite-field polynomials and conditionally linear samplers, and it checks
coefficient identities, degree accounts, dependency blocks, marginal laws,
and finite-instance histograms.  The operator-algebraic and compression
theorems remain imported.  This note makes
the description-level Turing-machine--lambda-calculus correspondence precise,
identifies the fixed point, and records the evidence boundary at every rung.
Part~I is a self-contained computability refresher for readers returning to
Turing machines and lambda calculus after a long interval; Part~II then puts
exact resource bounds on the same ideas.
\end{abstract}

\tableofcontents

\clearpage
\part{A refresher on computability, from the paper's viewpoint}
\label{part:refresher}
\setlength{\parskip}{0.30em}

The outermost layer of the argument is classical computation.  Before any
low-degree test or entangled strategy appears, the proof must be able to take
the finite text of one verifier, transform that text into another verifier,
and arrange for a verifier to use a description of itself.  This part rebuilds
that vocabulary from the ground up.  It is deliberately qualitative about
polynomial overheads; the exact costed statements begin in Part~II.

\section{Turing machines in the form used by the paper}
\label{sec:refresher-tm}

\subsection{Programs, configurations, and one step}

A Turing machine separates two things which ordinary mathematical notation
often blends together: a finite rule book and an unbounded but finitely used
memory.  For one tape, a convenient formal tuple is
\[
 M=(Q,\Sigma,\Gamma,\delta,q_0,q_{\mathrm h}),
\]
where \(Q\) is a finite set of control states, \(q_0\in Q\) is the initial
state, \(q_{\mathrm h}\in Q\) is the halting state,
\(\Sigma=\{0,1\}\) is the input alphabet,
\(\Gamma\supseteq\Sigma\cup\{\sqcup\}\) is a finite tape alphabet with blank
symbol \(\sqcup\notin\Sigma\), and
\[
 \delta:(Q\setminus\{q_{\mathrm h}\})\times\Gamma
 \longrightarrow Q\times\Gamma\times\{L,S,R\}
\]
is the transition function.  We take tapes to be indexed by \(\N\), so a
left move at cell zero leaves the head at zero.  The paper uses a multitape
version: a \(k\)-input machine has \(k\) input tapes, one work tape, and one
output tape.  Its transition reads all scanned symbols, changes the state,
writes on the work and output tapes, and moves each head by at most one cell.
The input tapes are read-only.  Adding separate accept and reject halting
states, or allowing writes on input tapes, changes none of the computability
or polynomial-overhead statements below.

A \emph{configuration} is a complete instantaneous description
\[
 C=(q,T_1,h_1,\ldots,T_{k+2},h_{k+2}):
\]
the current state, the contents \(T_i:\N\to\Gamma\) of every tape, and the
head locations \(h_i\in\N\).  Only finitely many cells are nonblank in every
reachable configuration.  The notation \(C\vdash_M C'\) means that one
application of \(\delta\) changes \(C\) into \(C'\).  Thus a computation is
not a metaphor: it is the uniquely determined path
\(C_0\vdash_M C_1\vdash_M C_2\vdash_M\cdots\) in the graph of configurations.

\paragraph{Why this matters here.}
Later, Cook--Levin replaces a bounded computation by local constraints on
adjacent configurations.  That replacement works because one step looks only
at the finite control and the cells under the heads.  The enormous object is
the history; the local dynamical law \(\delta\) is finite.  This is close to a
lattice model with a fixed local update rule, except that the update is
deterministic and the head positions are part of the state.

\paragraph{Worked one-step example.}
If \(\delta(q,0)=(q',1,R)\), then
\[
 (q,\ldots\sqcup\,\underline 0\,1\sqcup\ldots,h)
 \ \vdash_M\ 
 (q',\ldots\sqcup\,1\,\underline 1\sqcup\ldots,h+1).
\]
Only the state, the scanned cell, and the head location have changed.

The machine \emph{halts} on input \(x\) if some \(C_T\) has state
\(q_{\mathrm h}\).  Its output is the longest nonblank initial string on the
output tape.  We write \(M(x)\) for that string and \(M(x)=\bot\) if the path
never reaches \(q_{\mathrm h}\).  The instance running time is
\[
 \Time_{M,x}:=\min\{T:C_T\text{ has state }q_{\mathrm h}\},
\]
with value \(\infty\) if the set is empty.  We sometimes abbreviate this as
\(\Time_M(x)\).  Notice the distinction between \emph{has not halted yet} and
\emph{will never halt}; no finite observation of a run decides the latter in
general.

\paragraph{Why this matters here.}
The compression argument repeatedly says that a generated verifier finishes
within a supplied bound.  That is a statement about the number of transitions
on every relevant input, not merely about the mathematical function obtained
after ignoring time.  A timeout wrapper may safely reject after \(T\) steps,
because it is counting the very transitions used in the definition above.

\paragraph{Worked timing example.}
For the transition just displayed, if the new state \(q'\) is the halt state,
then the initial configuration is assigned time one.  A timeout of zero
rejects without applying the transition; a timeout of one obtains the output.

\subsection{A complete two-state equality machine}

Here is a small instance in the paper's multitape style.  The machine has two
one-bit input tapes \(A,B\), a work tape \(W\), and an output tape \(O\).  It
has two nonhalting states \(q_0,q_1\), plus the designated state
\(q_{\mathrm h}\).  A table entry lists
\((\hbox{next state};\hbox{new }W,\hbox{new }O;
  \hbox{moves of }A,B,W,O)\).
Rows are tested from top to bottom; they are disjoint except for the last
catch-all row, so the table specifies every possible scanned tuple.

\begin{center}
\small
\begin{tabularx}{\textwidth}{P{0.09\textwidth}P{0.24\textwidth}X}
\toprule
State & Scanned \((A,B,W,O)\) & Transition \\
\midrule
\(q_0\) & \((0,0,\sqcup,\sqcup)\) or
          \((1,1,\sqcup,\sqcup)\)
 & \((q_1;1,\sqcup;R,R,S,S)\) \\
\(q_0\) & \((0,1,\sqcup,\sqcup)\) or
          \((1,0,\sqcup,\sqcup)\)
 & \((q_1;0,\sqcup;R,R,S,S)\) \\
\(q_1\) & \((\sqcup,\sqcup,r,\sqcup)\), \(r\in\{0,1\}\)
 & \((q_0;r,r;S,S,S,S)\) \\
\(q_0\) & \((\sqcup,\sqcup,r,r)\), \(r\in\{0,1\}\)
 & \((q_{\mathrm h};r,r;S,S,S,S)\) \\
\(q_0,q_1\) & every tuple not covered above
 & \((q_{\mathrm h};0,0;S,S,S,S)\) \\
\bottomrule
\end{tabularx}
\end{center}

The final row makes the table total and rejects malformed inputs.  On legal
one-bit inputs the first step computes the equality bit on \(W\), the second
copies it to \(O\), and the third enters the halt state.  For \(A=B=1\), the
full three-step trace, showing the scanned cells, is
\[
\begin{array}{c|c|c|c|c|c|l}
i&q&A&B&W&O&\text{rule used}\\ \hline
0&q_0&\underline1&\underline1&\underline\sqcup&\underline\sqcup
  &\text{equal-input row}\\
1&q_1&1\underline\sqcup&1\underline\sqcup&\underline1&\underline\sqcup
  &\text{copy row with }r=1\\
2&q_0&1\underline\sqcup&1\underline\sqcup&\underline1&\underline1
  &\text{halt row with }r=1\\
3&q_{\mathrm h}&1\underline\sqcup&1\underline\sqcup&\underline1&\underline1
  &\text{halted.}
\end{array}
\]
Thus \(M_{=}(1,1)=1\) and \(\Time_{M_=,(1,1)}=3\).  The same trace with
work/output bit zero gives \(M_{=}(0,1)=0\).

\paragraph{Why this matters here.}
The example is intentionally mundane: a paper verifier is just a much larger
finite transition table whose output bit says accept or reject.  Statements
about a verifier's code, time, or simulation refer to this concrete object,
even when the paper describes it in higher-level pseudocode.

\subsection{Several inputs and the verifier convention}

Writing \(M(x_1,\ldots,x_k)\) means that \(x_i\) initially occupies the
initial cells of input tape \(i\), with all other cells blank and all heads at
zero.  Multiple tapes are convenient interfaces, not a stronger model: a
one-tape machine can store a self-delimiting encoding of all tapes and their
head markers and simulate one multitape step with polynomial overhead.

The paper's \emph{decider} is a total five-input machine
\[
 D(n,x,y,a,b)\in\{0,1\}.
\]
Here \(n\) is an index written in binary, \(x,y\) are the two questions, and
\(a,b\) are the two answers.  Total means that \(D\) halts on every such
five-tuple.  Its paper-level time at index \(n\) is the worst case over the
other four strings.  A normal-form verifier is a pair \(V=(S,D)\), where the
sampler \(S\) produces the questions and \(D\) checks the answers.

\paragraph{Why this matters here.}
The five inputs are not later bundled by sleight of hand.  When we translate a
decider to a lambda term, its argument is a five-component encoded tuple; when
we translate back, those components initialize five separate tapes.  The
worst-case convention is what lets a resource bound hold independently of
unread suffixes of arbitrarily long answer strings.

\paragraph{Worked verifier example.}
Let \(S\) always ask the one-bit questions \(x=y=0\), and let \(D\) ignore
\(n,x,y\) and run \(M_=\) on the first bits of \(a,b\).  Then the players are
accepted exactly when their one-bit answers agree, and the decision stage
takes three transitions plus fixed input-parsing overhead.

\subsection{Descriptions are finite data}

Because \(Q,\Gamma\), and \(\delta\) are finite, we may serialize them with a
fixed prefix-free convention.  The resulting bit string is the machine
description \(\langle M\rangle\); the paper also writes \(\overline M\).
Its length in bits is
\[
 |M|:=|\langle M\rangle|.
\]
Malformed strings are assigned a fixed rejecting machine, so every bit string
still has a defined interpretation.  The encoding is syntactic: extensionally
equivalent machines may have different codes, sizes, and running times.  A
useful physics analogy is that a description is like a written Hamiltonian,
not its spectrum.  The actual definition is the serialized state set and
transition table, and an algorithm may inspect those bits without solving the
machine's behavior.

\paragraph{Why this matters here.}
The size bound \(|V|\leq\lambda\) constrains the verifier's source code, not
the number of functions extensionally equivalent to it.  Transformations such
as introspection and answer reduction copy, scan, hardwire, and simulate that
source.  They cannot be defined merely on a black-box input/output relation.

\paragraph{Worked description example.}
The code of \(M_=\) contains tags for the three states and the five rows above.
A compiler can prepend a wrapper saying ``run this code, but reject after
three simulated steps.''  It produces new finite code without having to
decide any semantic property of \(M_=\).

A verifier description is the ordered pair
\(\langle V\rangle=(\langle S\rangle,\langle D\rangle)\).  In
\code{fig:compress}, \(\Compress\) literally receives
\((\langle V\rangle,\lambda)\), constructs descriptions for introspection,
answer reduction, and parallel repetition, and returns the description of a
new verifier.  This is why it is meaningful for \(\Compress\) to take ``a
verifier'' as input: it takes the finite string naming the verifier's rules,
not an infinite table of all possible plays.

\section{Universality: one machine can interpret all machines}
\label{sec:refresher-universality}

\subsection{The universal machine}

Fix an encoding of tuples and programs.  A universal machine \(U_k\) is a
single machine satisfying
\[
 U_k(\langle M\rangle,x_1,\ldots,x_k)=M(x_1,\ldots,x_k)
\]
whenever the right side halts, and diverging when that simulated run diverges.
Operationally, \(U_k\) keeps the code of \(M\) on a read-only region, stores an
encoded simulated configuration, finds the matching transition-table entry,
and updates the encoding.  Universality therefore says that descriptions can
be used as ordinary input data.

The paper needs the quantitative version.  With its dual-rail tuple encoding,
there are universal constants \(C_U,c_U\geq1\) such that a \(T\)-step run of a
\(k\)-input machine is reproduced within
\[
 C_U\bigl(k|\langle M\rangle||x|T\bigr)^{c_U},
 \qquad |x|=\sum_i|x_i|.
\]
No particular exponent is sacred.  What matters is that the overhead is one
fixed polynomial after the encodings have been fixed.

\paragraph{Why this matters here.}
Generated verifiers call old verifiers as subroutines.  Their finite source
contains the universal interpreter plus the old description; it does not
contain a separate hand-compiled copy of every possible transition.  The
polynomial overhead is later absorbed into a new resource exponent.

\paragraph{Worked universal simulation.}
On input \((\langle M_=\rangle,1,0)\), \(U_2\) parses the five table rows,
builds the initial four-tape configuration, and simulates the three transitions
above.  It outputs zero.  The interpreter is fixed; changing the two input
bits or replacing \(\langle M_=\rangle\) changes only its data.

\subsection{Parameter fixing: the
\texorpdfstring{\(s\)-\(m\)-\(n\)}{s-m-n} theorem}

The notation \(s\)-\(m\)-\(n\) packages a simple compiler operation.  If
program \(e\) expects \(r+k\) arguments, there is a total computable function
\(s_{r,k}\) on bit strings such that
\[
 \varphi^{(k)}_{s_{r,k}(e,z_1,\ldots,z_r)}(x_1,\ldots,x_k)
 =\varphi^{(r+k)}_e(z_1,\ldots,z_r,x_1,\ldots,x_k).
\]
The left program has the first \(r\) values hardwired into its source.

\paragraph{Proof idea in two lines.}
Print a wrapper containing one literal copy of \(e,z_1,\ldots,z_r\).  On input
\(x\), the wrapper forms the tuple
\((e,z_1,\ldots,z_r,x)\) and invokes the universal machine; printing this
wrapper is itself a terminating, effective source transformation.

\paragraph{Why this matters here.}
Hardwiring is the bridge from ``this value is supplied at run time'' to ``this
value is part of the verifier description.''  Compression constantly fixes a
verifier, a level, or a time exponent while leaving \((n,x,y,a,b)\) as the
live inputs.  Part~II calls the corresponding operation \(\Specialize\).

\paragraph{Worked specialization.}
Apply \(s_{1,1}\) to the two-input equality machine and the static bit \(1\).
The resulting one-input program has code
\(e_1=s_{1,1}(\langle M_=\rangle,1)\) and satisfies
\(\varphi_{e_1}(a)=M_=(1,a)\).  Thus it is the one-bit identity test, even
though its source is a universal-simulation wrapper rather than a newly
invented transition table.

\subsection{Why the Halting problem is undecidable}

Let
\[
 \mathsf{HALT}=\{(\langle M\rangle,x):M(x)\neq\bot\}.
\]
The Halting problem asks for a total decider \(H\) for this set.  The obstruction
is diagonalization, and it is worth seeing every step.

\begin{proof}[Diagonal proof]
Assume, for contradiction, that \(H(e,x)\) always halts and returns one exactly
when the program described by \(e\) halts on \(x\).  Construct a one-input
machine \(D\) as follows on input \(z\):
\begin{enumerate}
\item Run \(H(z,z)\).
\item If \(H\) returns zero, halt.
\item If \(H\) returns one, loop forever.
\end{enumerate}
This is a legitimate finite program because \(H\) was assumed to be a total
subroutine.  Let \(d=\langle D\rangle\) and run \(D(d)\).

If \(H(d,d)=0\), then by the specification of \(H\), \(D(d)\) does not halt;
but Step 2 says that it halts.  If \(H(d,d)=1\), then the specification says
that \(D(d)\) halts; but Step 3 says that it loops forever.  The two possible
output bits both yield contradictions.  Therefore no such total \(H\) exists.
\end{proof}

This proof does not say that a particular long run can never be understood.
It says there is no single terminating algorithm which answers the halting
question correctly for every program-description/input pair.

\paragraph{Why this matters here.}
The \(\mathrm{MIP}^*=\mathrm{RE}\) construction reduces the behavior of an
arbitrary machine to the value of a nonlocal game.  Its Halting verifier is
given \(\langle M\rangle\), embeds and simulates that description, and asks
whether bounded prefixes of the computation are consistent.  It does not
contain an impossible halting decider; rather, the infinite family of game
instances reflects the recursively enumerable process of eventually seeing a
halt.

\paragraph{Worked semidecision.}
For a fixed \(M,x\), simulate one more step at each stage \(t=0,1,2,\ldots\).
If \(M(x)\) halts after \(T\) steps, stage \(T\) witnesses it.  If it never
halts, this search itself never returns.  That one-sided behavior is the
``RE'' side of the theorem's computability interface.

\section{Programs which obtain their own descriptions}
\label{sec:refresher-kleene}

The diagonal argument used a program on its own code.  Kleene's second
recursion theorem turns that move into a general construction.  In its most
useful form here:

\begin{quote}
Every total computable transformation of programs has a program that behaves
like its own transform.
\end{quote}

More formally, if \(Q\) is a total computable map from descriptions of
\(k\)-input programs to descriptions of \(k\)-input programs, then there is a
description \(e_Q\) such that
\[
 \varphi^{(k)}_{e_Q}=\varphi^{(k)}_{Q(e_Q)}.
\]
This is behavioral equality, not necessarily literal equality of bit strings.

For the paper, this theorem buys a finite verifier description whose behavior
may depend on that very description.
\begin{theorem}[Kleene's second recursion theorem, refresher form]
\label{thm:refresher-kleene}
Every total computable transformation \(Q\) of \(k\)-input program codes has
an effectively constructible behavioral fixed point \(e_Q\).
\end{theorem}

\begin{proof}
Use parameter fixing for one static and \(k\) live arguments.  Write \(s(z,w)\)
for the code obtained by fixing the first input of a suitable \((k+1)\)-input
program \(z\) to the literal \(w\); hence
\[
 \varphi^{(k)}_{s(z,w)}(x)=\varphi^{(k+1)}_z(w,x).
\]
Now define the partial computable function
\[
 g(z,x):=\varphi^{(k)}_{Q(s(z,z))}(x).
\]
There is a finite \((k+1)\)-input program \(p\) for \(g\): on \((z,x)\), it
computes \(s(z,z)\), runs the total transformer \(Q\) on that string, and
universally simulates the resulting program on \(x\).  Set
\(e_Q:=s(p,p)\).  Then, for every \(x\),
\[
\begin{aligned}
 \varphi^{(k)}_{e_Q}(x)
 &=\varphi^{(k)}_{s(p,p)}(x)\\
 &=\varphi^{(k+1)}_p(p,x)\\
 &=\varphi^{(k)}_{Q(s(p,p))}(x)\\
 &=\varphi^{(k)}_{Q(e_Q)}(x).
\end{aligned}
\]
If the last computation diverges, each displayed partial computation diverges
for the same reason; if it halts, all return the same output.  Thus the
partial functions are equal on every input.  Every source-manipulating step
used to construct \(e_Q\) terminates, so the fixed point is effective.
\end{proof}

\paragraph{Why this matters here.}
The Halting verifier must pass its own decider description into \(\Compress\).
The theorem says that this is ordinary program construction, not a circular
definition requiring an infinite source string.  It also separates two facts:
constructing the fixed-point code terminates, while running the partial
program denoted by that code may or may not terminate.

\paragraph{Worked quine-like example.}
Let \(Q(z)\) be the code of the zero-input program ``print the literal string
\(z\).''  The transformation \(Q\) is total because it merely prints a
wrapper.  Its fixed point \(e_Q\) satisfies
\(\varphi_{e_Q}=\varphi_{Q(e_Q)}\), so running \(e_Q\) prints \(e_Q\) itself.
The theorem has built a quine without requiring the source printer to know its
own final address in advance.

\subsection{The paper's
\texorpdfstring{\(\mathcal F(\mathcal F,M,L,\ldots)\)}
{F(F,M,L,...)} construction}

The paper writes the fixed point in an unrolled form.  Begin with an
eight-input machine
\[
 \mathcal F(f,m,L,n,x,y,a,b).
\]
Its first action is a computable source transformation: hardwire
\((f,m,L)\) into the first three slots of \(\mathcal F\)'s template, obtaining
a five-input decider description \(d\).  It pairs \(d\) with a sampler
description, invokes \(\Compress\) on that verifier description, and runs the
resulting decider on \((n,x,y,a,b)\).  The published call is
\[
 \mathcal F(\langle\mathcal F\rangle,\langle M\rangle,L,n,x,y,a,b).
\]
After the first three arguments have been fixed, the produced \(d\) is exactly
the description of the five-input program currently being run.  Hence the
verifier passed to \(\Compress\) contains its own decider code.

\paragraph{Why this matters here.}
This is Kleene's construction with the diagonal substitution exposed as
pseudocode.  The transformation \(Q\) says ``insert the candidate code into a
verifier, compress it, then run the returned decider.''  The
\(\mathcal F(\mathcal F,\ldots)\) trick supplies \(s(p,p)\) directly.  An
earlier version of the paper invoked Kleene's recursion theorem explicitly;
the published presentation changes the notation, not the computability fact.

\paragraph{Worked schematic expansion.}
Suppress the fixed parameters \(M,L\) and denote the five live inputs by
\(u\).  For a proposed self-description \(f\), put
\(d_f=\operatorname{hardwire}(\langle\mathcal F\rangle,f)\).
Substituting \(f=\langle\mathcal F\rangle\) gives
\[
 d_{\langle\mathcal F\rangle}(u)
 =\mathcal F(\langle\mathcal F\rangle,u)
 =Q(d_{\langle\mathcal F\rangle})(u),
\]
which is the recursion-theorem equation with the compiler operation written
out rather than hidden under the symbol \(s\).

\section{Lambda calculus from scratch}
\label{sec:refresher-lambda}

Turing machines make state and local time explicit.  Lambda calculus makes
substitution and composition explicit.  We start with its untyped core,
because fixed-point combinators live there naturally, and only later return to
the typed, fuelled language used by the project.

\subsection{Terms, scope, and beta reduction}

A lambda term is generated by
\[
 t ::= x \mid (\lambda x.t) \mid (t\;t).
\]
A variable occurrence is \emph{bound} if it lies inside a matching
\(\lambda x\); otherwise it is \emph{free}.  For example, in
\(\lambda x.(x\;y)\), the occurrence of \(x\) is bound and \(y\) is free.
Changing a bound name consistently is \emph{alpha-renaming}:
\(\lambda x.x=_{\alpha}\lambda z.z\).  It changes typography, not the term's
binding structure.

Application associates to the left and abstraction extends to the right, so
\(f\;a\;b\) means \((f\;a)\;b\), while \(\lambda x.f\;x\) means
\(\lambda x.(f\;x)\).  The computational rule is beta reduction,
\[
 (\lambda x.t)\;u\longrightarrow_\beta t[x:=u],
\]
where substitution replaces the free occurrences of \(x\) in \(t\) and first
alpha-renames binders when necessary to avoid capture.

\paragraph{Why this matters here.}
Hardwiring a machine input and applying a function are both forms of
substitution.  Our later \(\Specialize\) operation is the source-code version:
it fills typed holes while preserving scope.  Being precise about bound and
free variables is what prevents a verifier description from changing meaning
when code is inserted under a binder.

\paragraph{Worked reduction 1.}
The identity combinator \(I=\lambda x.x\) satisfies
\[
 (\lambda x.x)\;((\lambda z.z)\;w)
 \longrightarrow_\beta (\lambda z.z)\;w
 \longrightarrow_\beta w.
\]

\paragraph{Worked reduction 2: avoiding capture.}
Naively substituting \(y\) for \(x\) in
\((\lambda x.\lambda y.x)\;y\) would incorrectly bind the free \(y\).
Rename the inner binder first:
\[
 (\lambda x.\lambda y.x)\;y
 =_\alpha(\lambda x.\lambda z.x)\;y
 \longrightarrow_\beta\lambda z.y.
\]
The surviving \(y\) is free, as it must be.

A \emph{beta-redex} is a subterm of the form \((\lambda x.t)u\).  A
\emph{normal form} contains no beta-redex.  Some terms have none: with
\(\Omega=(\lambda x.xx)(\lambda x.xx)\), one has
\(\Omega\to_\beta\Omega\) forever.  The Church--Rosser, or confluence,
theorem states that if one term reduces by different paths to \(u\) and \(v\),
then \(u\) and \(v\) have a common reduct.  In particular, a normal form, if
it exists, is unique up to alpha-renaming.  We use this theorem as orientation
and do not prove it here.

\paragraph{Why this matters here.}
Confluence says that the mathematical result does not depend on which redex
is chosen, provided reduction reaches a normal form.  Resource accounting is
stricter: different strategies may take very different numbers of steps, and
one may diverge while another reaches the answer.  The project therefore
fixes an evaluator instead of treating all beta paths as operationally equal.

\paragraph{Worked strategy example.}
For \(K=\lambda x.\lambda y.x\),
\[
 K\;I\;\Omega\longrightarrow_\beta I
\]
under normal order, which reduces the leftmost outermost redex first and never
evaluates the unused argument.  Call-by-value first tries to turn \(\Omega\)
into a value and therefore loops.  This difference is exactly why the usual
\(Y\) combinator needs an eta-delayed cousin under call-by-value.

\subsection{Church data: behavior as representation}

Pure lambda calculus has no primitive bits, tuples, or integers.  Church
encodings represent data by what eliminator it applies.  These encodings are
not the project's efficient runtime representation; they show that the bare
calculus is computationally expressive.

\paragraph{Booleans.}
Define
\[
 \mathsf{true}=\lambda t.\lambda f.t,
 \qquad \mathsf{false}=\lambda t.\lambda f.f.
\]
They choose the first or second continuation.  For example,
\(\mathsf{true}\;A\;B\to_\beta A\), whereas
\(\mathsf{false}\;A\;B\to_\beta B\).

\paragraph{Why booleans matter here.}
A verifier ultimately returns one acceptance bit.  In the resource language
we use a primitive bit sort, but its conditional has precisely the selection
behavior exhibited by these terms.

\paragraph{If--then--else.}
Define
\(\mathsf{if}=\lambda b.\lambda t.\lambda e.b\;t\;e\).  Then
\[
 \mathsf{if}\;\mathsf{false}\;A\;B
 \to_\beta\mathsf{false}\;A\;B
 \to_\beta B.
\]

\paragraph{Why the conditional matters here.}
Every decider is built from branches on parsed bits, timeout tests, and local
predicates.  Fixing their evaluation order makes ``the branch not taken''
operationally meaningful, especially when it contains a recursive call.

\paragraph{Pairs.}
Let
\[
 \mathsf{pair}=\lambda a.\lambda b.\lambda p.p\;a\;b,
 \quad \mathsf{fst}=\lambda p.p\;\mathsf{true},
 \quad \mathsf{snd}=\lambda p.p\;\mathsf{false}.
\]
Writing \(\langle A,B\rangle=\mathsf{pair}\;A\;B\),
\[
 \mathsf{fst}\;\langle A,B\rangle
 \to_\beta\mathsf{true}\;A\;B\to_\beta A.
\]

\paragraph{Why pairs matter here.}
Machine configurations, verifier descriptions, and sampler/decider outputs
are finite products.  The project gives tuples primitive storage so that
projection has an explicit small cost, but the lambda interface is the same.

\paragraph{Numerals.}
The Church numeral \(n\) is
\[
 \overline n=\lambda f.\lambda x.f^n x,
 \quad
 \overline0=\lambda f.\lambda x.x,
 \quad
 \overline2=\lambda f.\lambda x.f(fx).
\]
It represents iteration: \(\overline2\;s\;z\to_\beta s(sz)\).

\paragraph{Why numerals matter here.}
An index \(n\), a timeout, and a fuel counter are all iterated control data.
Church numerals prove representability, while binary primitive naturals in
Part~II prevent a unary blow-up in the resource bounds.

\paragraph{Successor.}
Define
\(\mathsf{succ}=\lambda n.\lambda f.\lambda x.f(nfx)\).  Then
\[
 \mathsf{succ}\;\overline2
 \to_\beta\lambda f.\lambda x.f(\overline2 f x)
 \to_\beta\lambda f.\lambda x.f(f(fx))=\overline3.
\]

\paragraph{Why successor matters here.}
It is the simplest example of a computable map on an encoding.  Later
compilers similarly build a new syntax tree around an old one and prove a
size increase, rather than treating the new program as an abstract function.

\paragraph{Addition.}
Define
\(\mathsf{add}=\lambda m.\lambda n.\lambda f.\lambda x.mf(nfx)\).  Thus
\[
 \mathsf{add}\;\overline1\;\overline2\;f\;x
 \to_\beta \overline1 f(\overline2 f x)
 \to_\beta f(f(fx)),
\]
so the result is \(\overline3\).

\paragraph{Why addition matters here.}
Fuel bounds and encoded addresses are assembled from smaller bounds.  This
example reminds us that arithmetic is itself computation; Part~II uses
efficient binary primitives and charges for their reference machines.

\paragraph{Predecessor by the pair trick.}
Predecessor is subtler because a Church numeral only says ``iterate.''  Let
\[
 \Phi=\lambda p.\mathsf{pair}\;(\mathsf{snd}\;p)
                 \;(\mathsf{succ}(\mathsf{snd}\;p)),
 \qquad
 \mathsf{pred}=\lambda n.\mathsf{fst}
   (n\;\Phi\;(\mathsf{pair}\;\overline0\;\overline0)).
\]
Each \(\Phi\)-step maps \(\langle j-1,j\rangle\) to
\(\langle j,j+1\rangle\).  The complete small calculation for three is
\[
 \langle0,0\rangle\xrightarrow{\Phi}\langle0,1\rangle
 \xrightarrow{\Phi}\langle1,2\rangle
 \xrightarrow{\Phi}\langle2,3\rangle
 \xrightarrow{\mathsf{fst}}2.
\]
Hence \(\mathsf{pred}\;\overline3\to_\beta^*\overline2\).

\paragraph{Why predecessor matters here.}
A timeout counts down, so total bounded simulation needs a predecessor and a
zero test even when the simulated program diverges.  The pair construction is
a miniature instance of adding auxiliary state to make a desired update
locally computable.

\paragraph{Lists.}
Define a fold encoding
\[
 \mathsf{nil}=\lambda c.\lambda z.z,
 \qquad
 \mathsf{cons}=\lambda h.\lambda t.\lambda c.\lambda z.
                    c\;h\;(t\;c\;z).
\]
Then the list \([A,B]=\mathsf{cons}\;A\;(\mathsf{cons}\;B\;\mathsf{nil})\)
satisfies
\[
 [A,B]\;C\;Z\longrightarrow_\beta^* C\;A\;(C\;B\;Z).
\]

\paragraph{Why lists matter here.}
Tapes, syntax children, environments, and continuation stacks are all finite
sequences with potentially unbounded length.  A functional list supplies the
shape; the resource IR adds heap pointers and charged operations so that the
cost of traversing it is visible.

\subsection{Recursion from self-application}

Suppose \(F\) maps a proposed recursive procedure to its body.  We want a term
\(f\) satisfying
\[
 f=Ff.
\]
Think of the occurrence of \(f\) on the right as a hole.  If the desired term
can be written as self-application \(h\;h\), then the equation asks for
\(h\;h=F(h\;h)\).  Choose
\[
 h=\lambda x.F(x\;x).
\]
Substituting this \(h\) into \(h\;h\) yields the fixed-point combinator
\[
 Y=\lambda F.(\lambda x.F(x\;x))(\lambda x.F(x\;x)).
\]
The promised equation is a direct beta calculation:
\[
\begin{aligned}
 Y\;F
 &\longrightarrow_\beta
 (\lambda x.F(xx))(\lambda x.F(xx))\\
 &\longrightarrow_\beta
 F\bigl((\lambda x.F(xx))(\lambda x.F(xx))\bigr)\\
 &=F(YF).
\end{aligned}
\]
There is no recursive name in the finite syntax of \(Y\); recursion appears
when two copies are applied.

\paragraph{Why this matters here.}
The lambda fixed point is the term-level analogue of Kleene's theorem.  Both
build a finite object whose unfolding supplies that object's own behavior.
Our IR will make the analogous code link explicit and constant-size, because
blindly expanding \(YF\) would obscure description length and fuel.

\paragraph{Worked fixed point.}
Take \(F=\lambda r.\lambda n.\mathsf{if}\;(\mathsf{iszero}\;n)
\;\overline1\;(\mathsf{mul}\;n\;(r(\mathsf{pred}\;n)))\).  Under normal-order
reduction, \(YF\;\overline3\) unfolds only when the selected branch needs it:
\[
\begin{aligned}
 YF\;3&\to F(YF)\;3\\
 &\to 3\cdot(YF\;2)
 \to 3\cdot2\cdot(YF\;1)\\
 &\to3\cdot2\cdot1\cdot(YF\;0)
 \to3\cdot2\cdot1\cdot1=6.
\end{aligned}
\]
Here the arithmetic names abbreviate their Church encodings; every displayed
arrow stands for their finite beta reductions.

Call-by-value evaluates an argument before entering the function body, so
\(YF\) tries to expand itself forever.  Eta-delay the recursive occurrence:
\[
 Z=\lambda F.
 (\lambda x.F(\lambda u.xxu))
 (\lambda x.F(\lambda u.xxu)).
\]
Now
\[
 ZF\to_\beta F(\lambda u.(ZF)u),
\]
and the recursive computation remains under a lambda until an actual argument
arrives.  Under call-by-value the conditional branches must likewise be
thunked.  With
\[
 F_v=\lambda r.\lambda n.
 \bigl(\mathsf{if}\;(\mathsf{iszero}\;n)\;
   (\lambda u.\overline1)\;
   (\lambda u.\mathsf{mul}\;n\;(r(\mathsf{pred}\;n)))\bigr)\;I,
\]
only the selected branch is applied to the dummy value \(I\), and
\(ZF_v\;\overline3\to_\beta^*\overline6\) under call-by-value.

\paragraph{Why evaluation strategy matters here.}
The project evaluator is deterministic call-by-value, so its fixed-point
constructor must behave like \(Z\), not raw \(Y\).  Once the strategy is
fixed, ``number of reductions'' becomes an honest resource that can be
translated to Turing-machine time.

\subsection{De Bruijn indices and quotation}

Names are pleasant for handwriting but awkward for a compiler: alpha-equal
terms have many spellings, and substitution must keep inventing fresh names.
De Bruijn notation replaces a bound occurrence by the number of binders
between it and its binder, counting the nearest binder as zero.  Thus
\[
 \lambda x.\lambda y.x\;y
 \quad\text{becomes}\quad
 \lambda.\lambda.(1\;0).
\]
Both \(\lambda a.\lambda b.a\;b\) and
\(\lambda x.\lambda y.x\;y\) serialize to the same tree.

\paragraph{Why this matters here.}
The IR needs canonical bytes and capture-free specialization.  De Bruijn
indices remove alpha-renaming from the representation: a variable is a
lexical address, and inserting a closed term cannot capture one of its free
names because it has none.

\paragraph{Worked index lookup.}
In \(\lambda.\lambda.(1\;0)\), applying the outer function to \(A\) and the
inner function to \(B\) creates environment frames \([B],[A]\).  Index zero
selects \(B\); index one crosses one frame and selects \(A\).  The body is
therefore \(A\;B\), just as the named term says.

To \emph{quote} a term is to treat its syntax rather than its value as data.
A term is a finite ordered tree: each node has a constructor tag
(variable, lambda, application, and so on), finite payloads, and finitely many
children.  Prefix-free integer and byte encodings serialize that tree to a
unique bit string \(\operatorname{ser}(t)\).  We define
\[
 |t|:=|\operatorname{ser}(t)|.
\]
Quotation does not claim to compute what \(t\) will do.  It merely reifies the
finite tree so an evaluator or compiler can inspect it.

\paragraph{Why this matters here.}
\(\Compress\) must accept code, measure its length, insert static parameters,
and return new code.  A runtime closure with a hidden host-language
environment is not adequate; a quoted, closed tree is.  Quotation is the
lambda-calculus version of writing \(\langle M\rangle\) on a universal
machine's input tape.

\paragraph{Worked quotation.}
The named identity terms \(\lambda x.x\) and \(\lambda z.z\) both become the
de Bruijn tree \(\mathsf{Lambda}(\mathsf{BoundVar}(0))\).  Quoting either
produces the same canonical bytes.  Evaluating those bytes on \(A\) returns
\(A\); copying the bytes returns code, not \(A\).

\section{Church--Turing as an engineering statement}
\label{sec:refresher-engineering}

The Church--Turing thesis is often remembered as a slogan about which
functions are computable.  Here we need a more prosaic engineering reading:
the two concrete programming media can compile into one another, and the
compilers preserve both finite descriptions and feasible bounds up to fixed
polynomials.

\paragraph{Turing machine to lambda term.}
Represent each tape as a functional zipper, represent a configuration as a
tuple, and write a fixed recursive loop which reads the hardwired transition
table and produces the next configuration.  One machine step becomes a
bounded collection of list, tuple, and table operations.  The term contains
one copy of \(\langle M\rangle\), so its serialized size is linear in
\(|M|\); a \(T\)-step machine run becomes a polynomial number of evaluator
steps.  Theorem~\ref{thm:tm-to-l} gives the formal bound.

\paragraph{Why this direction matters here.}
It shows that replacing a paper verifier by an \(\mathcal L\)-term loses none
of its behavior and makes no hidden exponential copy of its transition table.

\paragraph{Worked direction.}
For \(M_=\), the recursive loop starts from the four-tape tuple, looks up the
equal-input row, then the copy row, then the halt row.  Three loop iterations
produce the Church-level bit one, with fixed overhead for tuple and list
operations.

\paragraph{Lambda term to Turing machine.}
Write a fixed evaluator whose work tape stores the parsed syntax tree, an
environment, a continuation stack, and a fuel counter.  Each evaluation rule
is a finite transition routine.  A run using \(F\) charged reductions is
simulated in time polynomial in the code length, input length, and \(F\).
Hardwiring the quoted term into this evaluator gives a machine description
linear in the serialized term.  Theorem~\ref{thm:l-to-tm} makes this precise.

\paragraph{Why this direction matters here.}
The paper's theorems are stated for Turing-machine verifiers.  Compiling a
resource-accounted term back to that model is what licenses using the lambda
IR without silently changing the theorem being applied.

\paragraph{Worked direction.}
For the quoted term \((\lambda x.x)\;1\), the evaluator parses an application
node, builds the identity closure, places the bit in its environment, looks up
index zero, and emits one.  A Turing machine can represent every one of those
finite records on its work tape.

The promised dictionary is therefore:
\begin{center}
\small
\begin{tabularx}{0.94\textwidth}{P{0.29\textwidth}P{0.29\textwidth}X}
\toprule
Turing-machine notion & Lambda/IR notion & Operative correspondence \\
\midrule
Program & Closed term & Finite rule table versus finite syntax tree \\
Input tape contents & Argument term/value & Tuple components keep the inputs
separate \\
Machine transitions and time & Fuel/charged reduction steps & Each simulates
the other with fixed polynomial overhead \\
\(\langle M\rangle\), \(|M|\) & Serialized quoted tree, \(|t|\) & Canonical
finite data and its bit length \\
Universal machine & Evaluator & One fixed program interprets supplied code \\
\(s\)-\(m\)-\(n\) hardwiring & \(\Specialize\) & A computable map from code
and static data to new code \\
Kleene recursion theorem & \(Y\)/\(Z\), or the IR's \(\mathsf{Fix}\) & A
finite behavioral fixed point \\
\bottomrule
\end{tabularx}
\end{center}

\paragraph{What the dictionary does not say.}
It does not assert step-for-step equality, equal constants, or equal
description strings.  It asserts a pair of explicit compilers, semantic
agreement on each input, linear description growth, and polynomial running
time overhead.  Those are exactly the invariants needed when an exponent
\(\lambda\) is enlarged by a universal constant.

\section{Why this machinery appears in a paper about entangled provers}
\label{sec:refresher-why}

A nonlocal game is an analytic object: its value optimizes over strategies,
and in the entangled setting those strategies use states and measurements.
Nevertheless, the \emph{outer construction} of the game is algorithmic.  The
compression theorem starts from the finite description of a normal-form
verifier \(V=(S,D)\), together with a resource exponent \(\lambda\), and
returns another finite verifier description.  Its hypotheses constrain both
\(|V|\) and worst-case running times.  Consequently, extensional statements
such as ``\(D\) decides the same predicate'' are not enough: the proof must
know which code is embedded, how long it is, and how expensive its universal
simulation will be.

At the description level, the main pipeline is
\[
 \langle V\rangle
 \xrightarrow{\ \Introspect\ }
 \langle V^{(1)}\rangle
 \xrightarrow{\ \AnswerReduce\ }
 \langle V^{(2)}\rangle
 \xrightarrow{\ \Repeat\ }
 \langle V^{(3)}\rangle.
\]
Each arrow is a terminating computable map on finite strings.  Introspection
embeds a description so a verifier can check a smaller representation of its
own computation.  Answer reduction constructs a bounded-trace circuit and
then a PCP-style verifier.  Repetition constructs several coordinated copies
and an acceptance rule.  The composite is \(\Compress\).  Polynomial
bookkeeping proves that its output is still a legal bounded verifier.

\paragraph{Why the fixed point enters.}
For a supplied machine \(M\) and level parameter \(L\), the Halting verifier
must arrange that its decider is the one appearing inside the verifier fed to
\(\Compress\).  This is a fixed-point equation on descriptions.  Kleene's
theorem, the explicit \(\mathcal F(\mathcal F,M,L,\ldots)\) program, and the
lambda fixed-point constructor are three presentations of the same outer
operation.  The point is not mystical self-awareness: it is a finite compiler
which prints code containing an appropriate code literal.

\paragraph{Worked pipeline fragment.}
Suppose the input decider has code \(d\).  Answer reduction first constructs a
succinct circuit whose hardwired local transition rule is read from \(d\).
It then emits a new decider code \(d'\) that reconstructs and checks the
corresponding local constraints.  The transformation \(d\mapsto d'\) can be
run by a Turing machine, and its time and output length can be bounded before
either decider is executed.  Compression composes several transformations of
exactly this kind.

\paragraph{What is not computability theory.}
The dictionary stops at the boundary between constructing a game and proving
what quantum strategies can do in it.  The low-degree test's soundness,
rigidity of observables, state-dependent approximation estimates,
oracularization against entangled strategies, entanglement-dimension bounds,
and the quantitative preservation of completeness and soundness are analytic
statements.  They do not follow from universality, quotation, or a fixed-point
theorem.  A lambda implementation can expose the verifier transformation and
check finite algebraic identities; it cannot by that fact alone prove the
operator-algebraic leaves.

\paragraph{Why this separation matters here.}
It prevents two opposite mistakes.  One is to treat source manipulation as
informal and accidentally hide a noncomputable or super-polynomial step.  The
other is to imagine that a clean self-interpreter settles the quantum
soundness argument.  The paper needs both layers: computability theory with
polynomial bookkeeping on the outside, and the genuinely quantum analytic
theorems on the inside.  The resource-accounted description language in
Part~II is designed for the first layer and records exactly where it must call
the second.

\clearpage
\setlength{\parskip}{0.55em}
\part{Resource-accounted analytic underpinnings}
\label{part:accounted}

\section{The paper's machine model and its intensional obligations}
\label{sec:paper-machines}

\subsection{Machines, inputs, time, and descriptions}

We first repeat the conventions on which the later change of description
language depends.  This is necessary because the change is intensional: it
must preserve programs, bounds, and programs which inspect other programs,
not only partial functions.

\begin{reminder}
Section~\ref{sec:refresher-tm} distinguished the finite transition table from
its run.  In this part, \(\overline M\) is that table's canonical bit string,
\(|M|\) is its bit length, and \(\Time_{M,x}\) counts transitions of the run.
\end{reminder}

\begin{definition}[The paper's Turing-machine convention]
\label{def:paper-tm}
A \(k\)-input Turing machine has \(k\) input tapes, one work tape,
and one output tape.  Tapes are one-way infinite, indexed by \(\N\), and use
\(\{0,1,\sqcup\}\).  Work and output tapes are initially blank.  If the
machine halts, its output is the longest nonblank initial string on the output
tape, with the all-blank output interpreted as the one-bit string \(0\).  On
\(x=(x_1,\ldots,x_k)\), write \(M(x)\) for this string and write \(M(x)=\bot\)
if it never halts.  Its instance time
\(\Time_{M,x}\in\N\cup\{\infty\}\) is the number of transitions to the halt
state.

Every machine has canonical binary code \(\overline M\).  The code lists its
finite states and transition table using a fixed reasonable encoding; its bit
length is \(|\overline M|\), abbreviated \(|M|\).  For a bit string
\(\alpha\), \([\alpha]_k\) is the \(k\)-input machine described by \(\alpha\)
(with the fixed convention for malformed strings).  Tuples are encoded by
\(\operatorname{enc}_k\): each data bit is replaced by
\(0\mapsto01\), \(1\mapsto10\), and each component ends in \(00\).  Thus
\(|\operatorname{enc}_k(x)|=2\sum_i|x_i|+2k\), and encoding takes linear time.
\end{definition}

This is the model in the Turing-machine subsection of the paper.
\gt{sec:tms}{ground-truth/gt-03-prelim.tex}  Its universal-machine convention
is quantitative.  For each \(k\) there is a fixed one-tape \(U_k\), and
universal constants \(C_U,c_U\geq1\), such that if \([\alpha]_k(x)\) halts in
\(T\) steps then \(U_k(\operatorname{enc}_{k+1}(\alpha,x))\) produces the same
answer within
\[
 C_U\bigl(k|\alpha||x|T\bigr)^{c_U},
 \qquad |x|:=\sum_i|x_i|.
\]
The paper also assumes that a high-level machine which invokes a fixed
submachine has description length linear in the embedded submachine code.
Hardwiring a string \(z\) into one input of \(M\) is an effective source
transformation \(s(M,z)\); its output description has length
\(O(|M|+|z|+1)\), is produced in time \(O(|M|+|z|+1)\), and incurs polynomial
simulation overhead.  A timeout wrapper counts down in binary: before the
one-tape simulation it uses \(O(T\log T)\) time and \(O(|M|)\) description
length.  These are conventions, not exact equalities between particular
machine encodings.

\begin{definition}[Decider]
\label{def:paper-decider}
A decider is a five-input machine \(D\) which halts with one bit on every
\((n,x,y,a,b)\), where \(n\) is a positive integer in binary and
\(x,y,a,b\in\{0,1\}^*\).  The first input is the \emph{index}, and
\[
 \Time_D(n)=\max_{x,y,a,b}\Time_{D,(n,x,y,a,b)}.
\]
The maximum is understood in \(\N\cup\{\infty\}\); a finite bound means that
the machine cannot spend time reading arbitrarily remote symbols of the four
unrestricted input tapes.  Output \(1\) is acceptance and \(0\) is rejection.
\end{definition}

This is exactly \code{def:decider}.
\gt{def:decider}{ground-truth/gt-05-games-normalform.tex}

\begin{definition}[Conditionally linear sampler]
\label{def:paper-sampler}
A sampler \(S\) is a six-input Turing machine.  On index \(n\) its legal modes
are
\[
\begin{array}{ll}
(n,\mathsf{dimension}),&
(n,w,\mathsf{marginal},j,z),\\
(n,w,\mathsf{linear},j,u,y),&
(n,w,\mathsf{factor},j,u),
\end{array}
\]
where \(w\in\{\mathrm A,\mathrm B\}\).  Its outputs must describe the ambient
dimension, the \(j\)-th marginal, the conditional linear map, and the
coordinate indicator of its factor space for a pair of level-\(\ell\) CL
functions over \(\F_{q(n)}^{s(n)}\).  The paper writes \(\Time_S(n)\) for the
number of steps before \(S\) halts on index \(n\), i.e. the worst legal call
at that index.  Inputs and outputs which denote integers, finite-field
elements, vectors, and modes are represented by fixed binary conventions.
\end{definition}

This transcribes \code{def:sampler}; its four interfaces and the fact that no
call uses more than six tapes are explicit in the source.
\gt{def:sampler}{ground-truth/gt-04-cl.tex}

\begin{definition}[Normal form and \(\lambda\)-boundedness]
\label{def:paper-lambda}
A normal-form verifier is \(V=(S,D)\), where \(S\) is a CL sampler with field
size \(q(n)=2\) and \(D\) is a decider.  Its description length and level are
\[
 |V|=\max\{|S|,|D|\},\qquad
 \operatorname{level}(V)=\operatorname{level}(S).
\]
For an integer \(\lambda\geq1\), \(V\) is \(\lambda\)-bounded exactly when
\[
 |V|\leq\lambda,qquad
 \Time_S(n)\leq n^\lambda,qquad
 \Time_D(n)\leq n^\lambda\quad(n\geq2).
\]
No bound is required at \(n=1\).
\end{definition}

These are \code{def:normal-ver} and \code{def:lambda}.
\gt{def:normal-ver, def:lambda}{ground-truth/gt-05-games-normalform.tex}
In the game \(V_n\), question and answer alphabets are padded to lengths
\(\Time_S(n)\) and \(\Time_D(n)\), respectively.  The maximum-over-all-inputs
definition above is therefore stronger than a promise only about strings of
those lengths.

\subsection{What a verifier description means}

\begin{reminder}
A description is syntax as data, not an oracle for behavior.  The distinction
was illustrated by the code of the equality machine in
Section~\ref{sec:refresher-tm}; it is now used in anger because compilers scan,
copy, and hardwire verifier code.
\end{reminder}

A ``description of a verifier'' is the ordered pair
\((\overline S,\overline D)\) of canonical machine strings.  It is neither an
oracle for the functions computed by \(S,D\) nor an equivalence class of
programs.  An algorithm receiving it may count its bits, copy it, reject it as
malformed, put it on the input tape of a universal simulator, and manufacture
a new string with parts of it hardwired.  Two extensionally equal deciders can
have different \(|D|\), running times, and behavior under those source
transformations.

The compression theorem makes this interface literal.  There is a
polynomial-time machine \(\Compress\) whose input is
\((\overline V,\lambda)\), with \(\overline V=(\overline S,\overline D)\), and
whose output is the description of a nine-level verifier.  It calls, in order,
\begin{align*}
 \overline{V^{(1)}}&=\ComputeIntroVerifier(\overline V,\lambda,9),\\
 \overline{V^{(2)}}&=\ComputeAnsVerifier
      (\overline{V^{(1)}},\lambda,\mu,\gamma),\\
 \overline{V^{(3)}}&=\ComputeParrepVerifier
      (\overline{V^{(2)}},\lambda,\tau).
\end{align*}
Each call returns another pair of machine descriptions.
\gt{fig:compress, thm:compression}{ground-truth/gt-12-compression.tex}
Likewise, \code{thm:ar} states that \(\ComputeAnsVerifier\) ``outputs the
description of the answer reduced verifier,'' and its proof separately builds
the sampler and decider descriptions.
\gt{thm:ar}{ground-truth/gt-10-answer-reduction.tex}

\subsection{Universal and source-level obligations}

\begin{reminder}
Universality and \(s\)-\(m\)-\(n\), reviewed in
Section~\ref{sec:refresher-universality}, justify calling a supplied program
and fixing some of its arguments.  Every use below must additionally pay for
the universal simulation and for printing the specialized description.
\end{reminder}

The following list is the bookkeeping that a homoiconic description language
must discharge.  Calling the syntax ``lambda calculus'' removes none of these
obligations.

\begin{enumerate}
\item \textbf{Simulation and timeouts.}  Every high-level instruction which
calls \(S\), \(D\), or a generated verifier uses a universal machine, with the
polynomial overhead and timeout counter from \code{sec:tms}.  Hardwiring must
produce finite code with the asserted length.

\item \textbf{Introspection.}  The introspection decider contains the
description of the original sampler and invokes that sampler to obtain its
dimension, marginals, linear maps, and factor spaces.  The source constructs a
constant-size raw universal decider and then hardwires
\((\overline V,\lambda,\ell)\); it first scans \(|V|\) and replaces an
overlong description by a trivial verifier.  Thus introspection consumes the
sampler description, not merely its induced distribution.
\gt{lem:intro-decider-complexity, thm:introspection}{ground-truth/gt-08-introspection.tex}

\item \textbf{Answer reduction.}  In Step 1 of \code{fig:pcpverifier}, the
answer-reduced decider computes
\[
 \mathsf{Circuit}=\mathsf{PaddedSuccinctDecider}
   (\overline D,n,T,Q,\sigma,x,y)
\]
from the \emph{description} of \(D\).  The Cook--Levin circuit depends on the
transition table even though it is much smaller than the computation tableau.
The construction uses \(|D|\leq\sigma\) quantitatively.
\gt{fig:pcpverifier, prop:standard-succinct-sat}{ground-truth/gt-10-answer-reduction.tex}

\item \textbf{Compiler composition.}  The three machines displayed in the
preceding subsection traverse and assemble source descriptions.  The proof
that \(\Compress\) is polynomial time uses their running times.  The
independence of the compressed sampler is also a statement about which source
strings are embedded in its output.

\item \textbf{Self-reference.}  In the Halting section, the eight-input
machine \(\mathcal F\) is run as
\[
 \mathcal F(\overline{\mathcal F},\overline M,
             \lambda,n,x,y,a,b).
\]
It constructs the five-input program obtained by hardwiring its first three
arguments, pairs that decider description with a computed sampler description,
and passes the pair to \(\Compress\).  This is a source-level diagonalization,
not recursive invocation alone.
\gt{fig:halt_f}{ground-truth/gt-12-compression.tex}
\end{enumerate}

Sections~\ref{sec:resource-lambda}--\ref{sec:self-reference} construct one
description language which meets precisely these obligations.  The point is
not Church--Turing equivalence; the point is a costed equivalence between
finite, inspectable descriptions.

\section{A lambda calculus with resources}
\label{sec:resource-lambda}

\subsection{Syntax, phases, and canonical bytes}

\begin{reminder}
The bare calculus of Section~\ref{sec:refresher-lambda} had variables,
abstraction, application, beta reduction, and quotation.  The language
\(\mathcal L\) below is a typed, call-by-value engineering IR: de Bruijn
addresses make binding canonical, and explicit fuel makes evaluation cost a
first-class quantity.
\end{reminder}

Fix a finite set of ground sorts including \(\mathsf{Bit}\),
\(\mathsf{Nat}\), \(\mathsf{Bytes}\), \(\mathsf{Tuple}\),
\(\mathsf{MachineDesc}\), and \(\mathsf{Quoted}(A)\).  Function sorts are
finite products followed by a result sort.  The raw terms of the description
language \(\mathcal L\) are
\[
\begin{split}
t ::= {}& \operatorname{BoundVar}(d,j)
 \mid \operatorname{Hole}(h,A)
 \mid \operatorname{Lambda}(r,t)
 \mid \operatorname{Apply}(t,t^*)
 \mid \operatorname{Fix}(t)\\
&\mid \operatorname{If}(t,t,t)
 \mid \operatorname{Prim}(p,B,t^*)
 \mid \operatorname{Quote}(\operatorname{Closed}(t))\\
&\mid \operatorname{Eval}(t_{\rm code},t^*,F)
 \mid \operatorname{Specialize}(t_{\rm code},\rho).
\end{split}
\tag{8.1}\label{eq:l-grammar}
\]
Here \(t^*\) is a finite ordered list.  \(\operatorname{BoundVar}(d,j)\)
addresses slot \(j\) of the environment frame \(d\) binders outward; hence
alpha-conversion is absent from the representation.  A hole is a named,
typed, static substitution site, not a dynamic variable.  We impose the
\emph{affine-hole discipline}: each hole name occurs at most once.  Reuse of
static data is expressed by binding it once and using de Bruijn variables.
\(F\) is either \(\operatorname{FuelLiteral}(m)\) or a closed arithmetic term
\(\operatorname{FuelBound}(t_n,t_\lambda)\).

Each primitive name \(p\) has a total type contract, a deterministic reference
machine, and a natural-valued charge \(\kappa_p(v_1,\ldots,v_r)\).  The
serialized annotation \(B\) proves
\(\kappa_p(v)\leq B(|v_1|,\ldots,|v_r|)\).  The registry contains the usual
finite byte, bit, tuple, list, comparison, and arithmetic operations, as well
as the bounded machine-table operation used in
Theorem~\ref{thm:tm-to-l}.  A primitive is not an oracle: its reference
machine runs in at most
\(c_p(1+\kappa_p+\sum_i|v_i|)^{e_p}\) Turing steps, with \(c_p,e_p\) fixed in
the language definition.

The sorting judgment \(\Gamma;H\vdash t:A\) is the ordinary simply typed
call-by-value judgment for variables, abstraction, application, and
conditionals, extended as follows.  \(H\) records affine holes;
\(\Quote\) requires a scoped closed term with empty \(H\);
\(\Specialize\) requires an environment covering exactly \(H\) with terms of
the declared sorts; \(\Eval\) requires quoted code, arguments matching the
quoted function sort, and natural fuel.  \(\operatorname{Fix}(P)\) requires a partial body
\(P:A\) with the single hole
\(\mathsf{self\_code}:\mathsf{Quoted}(A)\), and closes that hole.  We write
\(P\{A\}\) for sorted syntax, \(\operatorname{Closed}(P)\) for empty variable
and hole contexts, and \(\operatorname{Quoted}\{A\}\) for its canonical bytes.
Raw ill-sorted syntax is retained so that the evaluator can return
\(\mathsf{TypeError}\) rather than throw a host exception.

\begin{definition}[Canonical serialization]
\label{def:l-serialization}
Let \(\nu(m)=1^\ell0b\), where \(b\) is the \(\ell\)-bit binary expansion of
\(m+1\) and \(\ell=\lfloor\log_2(m+1)\rfloor+1\).  Thus
\(|\nu(m)|=2\ell+1\), and \(\nu\) is prefix-free.  Encode a byte string \(z\)
as \(\nu(|z|)z\).  Encode a term by a four-bit constructor tag followed, in
the order displayed in \eqref{eq:l-grammar}, by \(\nu\)-encoded integer and
string fields, the encoded child count, and the recursive child encodings.
Sorts, primitive contracts, fuel expressions, and static environments use the
same rule with fixed tags and lexicographically ordered field names.  No node
contains a redundant subtree length.

The resulting map \(\operatorname{ser}\) is injective and self-delimiting.
For a term or code value \(t\), define \(|t|:=|\operatorname{ser}(t)|\) in
bits.  This is the description length used below.
\end{definition}

Prefix-freeness of \(\nu\) determines the end of every field; arities determine
the number of recursive parses.  Induction on the first constructor tag
therefore gives a unique parse.  In particular, description size is a number
computed from syntax, not the size of a runtime closure.

The operative representations are:
\begin{center}
\begin{tabular}{P{0.19\textwidth}P{0.34\textwidth}P{0.35\textwidth}}
\toprule
Object & Contents & Legitimate consumer \\
\midrule
Closure & Body pointer and runtime environment & Evaluator only; it has no
canonical project byte string \\
Quoted syntax & Canonical bytes of closed \(\mathcal L\) syntax & Evaluator,
specializer, compiler, and \(\Compress\) \\
Circuit & Finite Boolean DAG with named input blocks & Succinct clause
evaluator and Tseitin transformation \\
Universal evaluator & One fixed program interpreting quoted syntax & Timeout
and bounded-trace construction \\
Specialization & Capture-free code-to-code replacement & Source
transformations; its result is syntax, not a closure \\
\bottomrule
\end{tabular}
\end{center}

\subsection{Small-step semantics and fuel}

\begin{reminder}
Normal form alone does not specify the cost of reaching it.  Recall the term
\(K I\Omega\): normal order terminates and call-by-value does not.  We now fix
one deterministic call-by-value transition system, so a fuel count denotes a
particular run rather than an arbitrary beta-reduction path.
\end{reminder}

We use a deterministic call-by-value CEK machine.  A runtime closure is
\([\lambda r.t,\eta]\), where \(\eta=(\eta_0,\eta_1,\ldots)\) is a stack of
frames and \(\eta_d[j]\) supplies \(\operatorname{BoundVar}(d,j)\).  Other
values are canonical ground data and \(\operatorname{Code}(t)\) for closed
syntax.  A configuration is
\[
 C=\langle q,\eta,K,H,f\rangle,
\]
where \(q\) is the current term or value, \(K\) is a continuation stack,
\(H\) is an append-only heap for lists, tuples, closures, and syntax pointers,
and \(f\in\N\) is remaining fuel.  Arguments are evaluated left to right.

For reference, the contractions after all indicated subterms have become
values are the following.  The notation \(C\longrightarrow_k C'\) means that
the transition has charge \(k\); continuations and heaps not shown are carried
unchanged.
\[
\begin{array}{rcll}
\langle\operatorname{BoundVar}(d,j),\eta,K,H,f\rangle
 &\longrightarrow_1&\langle\eta_d[j],\eta,K,H,f-1\rangle
 &\text{if the address exists},\\
\langle\operatorname{Lambda}(r,t),\eta,K,H,f\rangle
 &\longrightarrow_1&\langle[\lambda r.t,\eta],\eta,K,H,f-1\rangle,&\\
\langle\operatorname{Apply}([\lambda r.t,\eta],\vec v),K,H,f\rangle
 &\longrightarrow_1&\langle t,(\vec v)::\eta,K,H,f-1\rangle
 &(|\vec v|=r),\\
\langle\operatorname{If}(1,v_1,v_0),K,H,f\rangle
 &\longrightarrow_1&\langle v_1,K,H,f-1\rangle,&\\
\langle\operatorname{If}(0,v_1,v_0),K,H,f\rangle
 &\longrightarrow_1&\langle v_0,K,H,f-1\rangle,&\\
\langle\operatorname{Prim}(p,B,\vec v),K,H,f\rangle
 &\longrightarrow_{\kappa_p(\vec v)}&
 \langle p(\vec v),K,H,f-\kappa_p(\vec v)\rangle,&\\
\langle\operatorname{Quote}(\operatorname{Closed}(t)),K,H,f\rangle
 &\longrightarrow_1&\langle\operatorname{Code}(t),K,H,f-1\rangle.&
\end{array}
\tag{8.2}\label{eq:cek-core}
\]
An absent variable address, failed primitive contract, nonbit condition,
wrong application arity, or noncode argument to Eval/Specialize has the sole
successor \(\mathsf{TypeError}\) at charge one.  Evaluation contexts supply
the omitted rules: for example, Apply first pushes a frame remembering its
unevaluated arguments, evaluates the operator, then evaluates each argument
from left to right.  Thus the display specifies contractions while the context
discipline specifies their unique order.

All administrative CEK transitions cost one unit.  They include variable
lookup, creating a closure, pushing or popping one continuation frame,
selecting a Boolean branch, and returning a value.  Beta reduction of
\(\operatorname{Apply}([\lambda r.t,\eta],v_1,\ldots,v_r)\) costs one and
continues with \(t\) in environment
\(((v_1,\ldots,v_r),\eta_0,\eta_1,\ldots)\).  A wrong arity, applying a
nonclosure, or using a nonbit condition takes one transition to
\(\mathsf{TypeError}\).  A residual \(\operatorname{Hole}\) does the same.
\(\operatorname{Quote}(\operatorname{Closed}(t))\) takes one transition to the
syntax pointer \(\operatorname{Code}(t)\); serialization was performed when
the enclosing description was admitted, so this transition does not copy
\(|t|\) bits.

After its arguments have become values,
\(\operatorname{Prim}(p,B,v_1,\ldots,v_r)\) costs exactly
\(\kappa_p(v_1,\ldots,v_r)\geq1\).  If the contract or dynamic sorts fail it
uses one unit and returns \(\mathsf{TypeError}\).  For purposes of the local
semantics, a charge \(k\) is expanded into \(k\) deterministic microsteps of
the primitive reference machine; thus every microstep consumes one fuel unit.

\(\operatorname{Specialize}(\operatorname{Code}(P),\rho)\) first checks the
sort and coverage of \(\rho\), then traverses \(P\) in prefix order.  It uses
\(1+|P|+\sum_h|\rho(h)|\) fuel and returns the code value of the substituted
term, or \(\mathsf{TypeError}\).  \(\operatorname{Eval}\) evaluates its code,
arguments, and fuel expression, installs a delimiter holding the requested
inner budget, and runs the same CEK transition relation inside it.  Installing
and removing that delimiter costs two units.  Each inner microstep decrements
both the inner budget and the ambient budget.  Exhausting either produces
\(\mathsf{OutOfFuel}\).

Finally, \(\operatorname{Fix}(P)\) is a value of the result sort of \(P\).
When it is applied to arguments \(u\), one unfolding transition installs the
static binding
\[
 \mathsf{self\_code}\longmapsto
 \operatorname{Code}(\operatorname{Fix}(P))
\]
and evaluates \(P\) with the same arguments.  This is a constant-size heap
link, not a serialization of a recursively expanded tree.  Including the
application and binding frames, an unfolding costs the fixed constant
\(c_Y=3\).  The corresponding fully specialized syntax can be materialized by
\(\Specialize\) when a source consumer asks for it.

If a transition has charge \(k\) and \(f<k\), its unique successor is
\(\mathsf{OutOfFuel}\); otherwise the successor has counter \(f-k\).  A final
configuration is exactly one of
\[
 \mathsf{Value}(v),\qquad \mathsf{OutOfFuel},\qquad
 \mathsf{TypeError}.
\]
We write \(\operatorname{run}(t;f)\) for this result.  For closed function
syntax \(t\) and an argument tuple \(u\),
\(\operatorname{eval}_{\mathcal L}(t,u;f)\) abbreviates the run of
\(\operatorname{Apply}(t,\overline u)\), where \(\overline u\) is the literal
value syntax.  Results do not record unused fuel.

For the paper, the next theorem buys a unique bounded-evaluation trace for
each quoted verifier call, which can later be encoded by Cook--Levin.
\begin{theorem}[Determinism and outcome trichotomy]
\label{thm:l-determinism}
Every nonfinal configuration with positive fuel has at most one successor.
For every closed raw term and finite initial fuel, evaluation terminates in
exactly one of \(\mathsf{Value}(v)\), \(\mathsf{OutOfFuel}\), or
\(\mathsf{TypeError}\).
\end{theorem}

\begin{proof}
The current control item is either a term, a value, or a primitive microstate.
Constructor tags are disjoint.  For a term tag, the evaluation-context rule
selects the leftmost child not yet represented by a value; if there is one,
exactly one continuation frame is pushed.  If all children are values, exactly
one contraction above applies.  Variable addresses select at most one frame
and slot.  Primitive reference machines and their contracts are deterministic
by definition.  The Boolean and arity tests have disjoint success and error
cases.  The same reasoning applies to the unique top Eval delimiter and to the
single self-code binding of Fix.  Hence two distinct successors cannot exist.

Every nonfinal transition consumes at least one unit.  Starting from fuel
\(f\), at most \(f\) transitions or primitive microsteps can occur before a
value or error is returned; if neither has occurred, the next requested step
returns out-of-fuel.  The three final tags are disjoint, and determinism makes
the reached tag unique.
\end{proof}

For the paper, the next theorem buys safe enlargement of a proved timeout:
extra fuel cannot change a completed verifier answer.
\begin{theorem}[Fuel monotonicity]
\label{thm:fuel-monotonicity}
If \(\operatorname{run}(t;f)=R\) with
\(R\neq\mathsf{OutOfFuel}\), then
\(\operatorname{run}(t;f+g)=R\) for every \(g\geq0\).  More precisely, if the
first run uses \(c\leq f\) units, the second uses the same \(c\) transitions
and leaves \(g+(f-c)\) units.
\end{theorem}

\begin{proof}
Induct on the \(c\) transitions of the terminating run.  At step \(i\), the
larger-fuel run has exactly \(g\) more ambient fuel and, inside each Eval
delimiter, exactly the same inner budget as the smaller run.  The smaller run
can pay the uniquely determined charge \(k_i\), so the larger run can pay it
as well.  Determinism gives the same nonfuel successor data, and subtraction
preserves the difference \(g\).  After \(c\) steps the larger run therefore
reaches the same final tag and value.  No operational rule branches on the
numeric fuel except the insufficient-fuel rule, which was not selected in the
smaller run.
\end{proof}

\subsection{Quotation, evaluation, and specialization}

\begin{reminder}
Quotation turns a finite syntax tree into data; specialization is the lambda
counterpart of \(s\)-\(m\)-\(n\).  Neither operation decides what the quoted
program computes.  It only parses or rewrites its canonical tree.
\end{reminder}

External evaluation of bytes must pay for decoding even though an internal
Quote is a syntax pointer.  Define
\[
 h(d,u)=3+|d|+|\operatorname{enc}(u)|.
\]
The quoted evaluator \(\Eval_{\mathcal L}(d,u;F)\) checks and decodes \(d\),
installs the argument tuple, and then runs the CEK machine.  By definition the
front-end scan uses exactly \(h(d,u)\) fuel; malformed code returns
\(\mathsf{TypeError}\) during that scan.

For the paper, the next theorem buys the right to replace execution of a term
by execution of its transmitted description while charging for decoding.
\begin{theorem}[Quote/Eval equation with explicit overhead]
\label{thm:quote-eval}
Let \(t\) be a closed function term, let \(d=\operatorname{ser}(t)\), and let
\(u\) have the required argument sorts.  For every \(f\geq0\),
\[
 \Eval_{\mathcal L}(\Quote(t),u;f+h(d,u))
   =\operatorname{eval}_{\mathcal L}(t,u;f).
\tag{8.3}\label{eq:quote-eval}
\]
Here the left side takes the canonical bytes denoted by \(\Quote(t)\).  If the
right side uses \(c\leq f\) units and returns a value or type error, the left
side uses exactly \(h(d,u)+c\).  If it requests a \((f+1)\)-st unit, both sides
return out-of-fuel after their respective budgets are exhausted.
\end{theorem}

\begin{proof}
The prefix decoder is inverse to serialization by
Definition~\ref{def:l-serialization}; hence after its fixed scan the left
machine holds the same syntax tree \(t\), literal arguments, empty lexical
environment, and application continuation as the right machine.  The left
counter is then exactly \(f\).  From this point the two configurations are
identical.  Theorem~\ref{thm:l-determinism} gives identical successors until a
final outcome.  The front-end accounts for the stated additional
\(h(d,u)\) units and no other rule differs.
\end{proof}

\begin{lemma}[Capture-free specialization and its size]
\label{lem:specialization}
Let \(P\) be well-scoped partial syntax satisfying the affine-hole discipline,
and let \(\sigma\) map every hole of \(P\) to closed, sort-correct syntax.
There is a unique term \(\Specialize(P,\sigma)\), it has no holes, it preserves
the sort of \(P\), and
\[
 |\Specialize(P,\sigma)|
 \leq |P|+\sum_{h\in\operatorname{dom}\sigma}|\sigma(h)|.
\tag{8.4}\label{eq:specialize-size}
\]
The construction takes at most \(1+|P|+\sum_h|\sigma(h)|\) CEK microsteps.
\end{lemma}

\begin{proof}
Traverse \(P\) in prefix order.  Copy every nonhole tag and payload unchanged.
At \(\operatorname{Hole}(h,A)\), check the unique table entry, check its sort,
and emit \(\sigma(h)\).  The inserted term is closed, so its de Bruijn indices
refer to no binder in \(P\); indices in the surrounding syntax are also
unchanged.  Thus no variable can be captured.  Induction on \(P\) proves sort
preservation, because the hole and replacement have the same sort at the only
new case.  Coverage removes every hole, and deterministic traversal gives
uniqueness.

Serialization has no ancestor subtree-length fields.  Replacing the one
occurrence of \(h\) deletes a nonempty encoded Hole node and inserts exactly
\(|\sigma(h)|\) bits.  Since distinct affine holes have disjoint occurrences,
\[
 |\Specialize(P,\sigma)|
 =|P|-\sum_h|\operatorname{Hole}(h,A_h)|+\sum_h|\sigma(h)|,
\]
which implies \eqref{eq:specialize-size}.  One scan step per input or emitted
bit, plus the initial check, gives the time bound.  Without the affine-hole
discipline the correct last term would sum once per hole occurrence; this is
why the discipline is part of the formal system.
\end{proof}

\section{Simulation theorems: the machine--term bridge}
\label{sec:simulation}

\begin{reminder}
Section~\ref{sec:refresher-engineering} gave the two compiler ideas and warned
that Church--Turing equivalence alone says nothing quantitative.  This section
supplies the description-length and running-time inequalities needed by the
compression theorem.
\end{reminder}

\subsection{From the paper's machines to
\texorpdfstring{\(\mathcal L\)}{L}}

We use explicit functional tapes and one bounded table primitive.  A tape is a
zipper \((L,s,R)\): \(L\) is the reversed list strictly left of the head,
\(s\in\{0,1,\sqcup\}\) is the scanned symbol, and \(R\) is the list strictly
right of it with implicit blank tail.  A right move maps
\((L,s,rR')\) to \((sL,r,R')\), taking \(r=\sqcup\) when \(R\) is empty; a
left move is symmetric.  These are fixed list primitives of charge at most
four.  A configuration is a tuple of the finite control state and the
\(k+2\) zippers for input, work, and output tapes.

For a machine description \(m\), the primitive
\(\mathsf{TMDelta}(m,q,s_1,\ldots,s_{k+2})\) scans the encoded transition table
and returns the next state, written symbols, and head motions.  Its declared
charge is
\[
 \kappa_{\mathsf{TMDelta}}\leq 8(|m|+k+2).
\]
The primitive validates \(m\); malformed descriptions use a fixed rejecting
transition.  Its implementation is an ordinary finite table scan.  Thus the
description of the table remains data of length \(|m|\), rather than being
duplicated into a nest of conditionals.

Let \(\mathsf{init}_k\) parse the \(k\)-tuple encoding into input zippers and
blank work/output zippers, and let \(\mathsf{out}\) read the longest nonblank
output prefix.  Define
\[
\begin{split}
 \mathsf{loop}_m={}&\operatorname{Fix}\bigl(
  \lambda C.\ \operatorname{If}(\mathsf{halt}(C),\mathsf{out}(C),\\[-2pt]
 &\hspace{8em}\mathsf{self}(\mathsf{move}(
       C,\mathsf{TMDelta}(m,\mathsf{view}(C)))))\bigr),\\
 \langle\!\langle M\rangle\!\rangle={}&\lambda z.\mathsf{loop}_{\overline M}
                  (\mathsf{init}_k(z)).
\end{split}
\tag{9.1}\label{eq:tm-interpreter-term}
\]
Here \(\mathsf{self}\) is the closure obtained from the Fix self-code link;
equivalently it is \(\Eval\) of that link with the next configuration.  All
other displayed operations are fixed primitives with the stated list
representation.

For the paper, the next theorem buys a lambda description for every original
sampler or decider without changing its answers or hiding a size blow-up.
\begin{theorem}[TM to \(\mathcal L\), with resources]
\label{thm:tm-to-l}
There are universal constants \(a_0,b_0,c_0\geq1\) such that for every
\(k\)-input machine \(M\), \eqref{eq:tm-interpreter-term} is a closed term and
\[
 |\langle\!\langle M\rangle\!\rangle|\leq a_0|M|+b_0.
\]
If \(M(x_1,\ldots,x_k)\) halts within \(T\geq1\) transitions, then
\[
 \Eval_{\mathcal L}\!\left(
   \Quote(\langle\!\langle M\rangle\!\rangle),
   \langle x_1,\ldots,x_k\rangle;
   c_M T(T+|x|)\right)=M(x_1,\ldots,x_k),
\tag{9.2}\label{eq:tm-to-l-bound}
\]
where \(|x|=\sum_i|x_i|\) and one may take
\(c_M=c_0(|M|+k+1)\).  The tuple on the left is the same dual-rail encoding as
in Definition~\ref{def:paper-tm}.  In particular, a decider receives exactly
\[
 \langle\operatorname{bin}(n),x,y,a,b\rangle_{\mathcal L},
 \qquad
 \operatorname{enc}_{\mathcal L}\bigl(\langle\operatorname{bin}(n),x,y,a,b
   \rangle_{\mathcal L}\bigr)
 =\operatorname{enc}_5(\operatorname{bin}(n),x,y,a,b),
\]
and the five components initialize five separate input zippers.
\end{theorem}

\begin{proof}
The term contains one copy of \(\overline M\), the fixed interpreter body, and
self-delimiting constructor overhead.  Each input bit in \(\overline M\)
contributes a fixed number of serialized bits, so constants \(a_0,b_0\)
independent of \(M\) give the size bound.

Initialization scans the dual-rail input once.  For valid input it consumes at
most \(c_0(|x|+k+1)\) units and produces exactly the paper's initial tapes.
Assume inductively that the functional configuration before loop iteration
\(i\) represents the machine configuration after \(i\) transitions.  The
views return the same scanned symbols.  The table primitive selects the unique
transition of \(M\), and each zipper update writes the selected symbol and
moves in the selected direction.  Therefore the next functional configuration
represents the machine configuration after \(i+1\) transitions.  This proves
the simulation invariant for \(0\leq i\leq T\).

One loop iteration, including the fixed-point unfolding, table scan, tuple
operations, and \(k+2\) zipper updates, uses at most
\(c_0(|M|+k+1)\) units.  No visited zipper has more than \(|x|+T\) explicit
cells.  The final output scan therefore uses at most
\(c_0(|x|+T+1)\) units.  Adding initialization, \(T\) iterations, output, and
the quote/decode overhead from Theorem~\ref{thm:quote-eval} gives at most
\[
 c_0(|M|+k+1)\bigl(|x|+1+T\bigr)+
 c_0(|M|+k+1)T.
\]
For \(T\geq1\), both \(|x|+T+1\) and \(T\) are at most
\(2T(T+|x|)\).  Increasing the universal choice of \(c_0\) by a factor four
makes this no greater than the fuel in \eqref{eq:tm-to-l-bound}.  The invariant
and the definition of \(\mathsf{out}\) then give the same output string.  The
explicit quadratic form is deliberately conservative; the zipper simulation
itself is linear in \(T\) apart from the hardwired table scan.
\end{proof}

\subsection{From \texorpdfstring{\(\mathcal L\)}{L} to the paper's machines}

For the paper, the next theorem buys a route back from the IR to the exact
Turing-machine model in which verifier bounds and compression are stated.
\begin{theorem}[A universal evaluator for \(\mathcal L\)]
\label{thm:l-to-tm}
There is a fixed three-input Turing machine \(U_{\mathcal L}\), of description
length \(b_U=O(1)\), and a universal constant \(C_L\geq1\), such that on input
\((d,u,f)\) it returns the same value, out-of-fuel, or type-error result as
\(\Eval_{\mathcal L}(d,u;f)\), and
\[
 \Time_{U_{\mathcal L}}(d,u,f)
 \leq C_L\bigl(|d|+|u|+f+1\bigr)^6.
\tag{9.3}\label{eq:l-to-tm-bound}
\]
Given closed \(t\), hardwiring \(d=\operatorname{ser}(t)\) produces a machine
\((U_{\mathcal L})_t\) with
\[
 |(U_{\mathcal L})_t|\leq a_U|t|+b_U.
\]
On the paper's separate input tapes it may represent each argument by a lazy
head-and-position handle.  Consequently, a run with fuel \(f\) accesses at
most \(f\) symbols of those tapes and has the refined bound
\(C_L(|t|+f+1)^6\), independent of unread suffixes.  Hence a total decider term
with fuel schedule \(f_D(n)\) is a decider in the paper's sense and
\[
 \Time_D(n)\leq C_L(|t|+f_D(n)+1)^6.
\tag{9.4}\label{eq:decider-term-time}
\]
\end{theorem}

\begin{proof}
The machine first checks the prefix serialization and writes node records on a
work-tape heap.  It stores a node address, environment-frame addresses,
continuation records, and binary fuel counters.  During at most \(f\)
microsteps the heap acquires at most a fixed number of records per step, so its
length is at most \(c(|d|+|u|+f+1)\) for a universal \(c\).  A conservative
single-tape implementation finds a record by a full scan, performs binary
address arithmetic, and returns to the simulated head in at most the square of
that length.  Decoding, \(f\) simulated steps, and final serialization are
therefore bounded by a universal constant times the fourth power of the
length.  Converting the fixed multitape implementation to the paper's
one-tape convention and allowing the paper's polynomial universal-simulation
overhead is bounded by the sixth power after increasing \(C_L\).  This proves
\eqref{eq:l-to-tm-bound}.  Primitive microstates use their fixed reference
machines and have charged running time; their polynomial exponents were fixed
when \(\mathcal L\) was fixed and are included in the exponent six (or, for a
different registry, in another fixed exponent).

The simulator description contains only the parser, CEK transition table,
primitive registry, and serializers, so its length is \(b_U\), independent of
\(d,u,f\).  The paper's hardwiring convention embeds \(d\) once and gives the
linear description bound.  For separate input tapes, a handle stores a tape
number and head position.  One microstep can move such a head by at most the
work performed during that charged step, so no more than \(f\) input symbols
can be inspected.  Removing the untouched encoded suffixes from the preceding
storage account yields the refined bound.  Taking a maximum over all
\((x,y,a,b)\) now proves \eqref{eq:decider-term-time}; arbitrarily long unread
suffixes do not affect it.
\end{proof}

\subsection{The resource dictionary}

\begin{reminder}
The qualitative dictionary in Section~\ref{sec:refresher-engineering} matched
machine time with reduction fuel and machine code with serialized syntax.  The
table below is the quantitative version: it records the polynomials rather
than merely asserting that simulations exist.
\end{reminder}

For a term program, let \(s_t=|t|\) and let \(f_t(n)\) be a proved fuel
schedule which suffices for every legal call at index \(n\).  The translations
above give the following dictionary.

\begin{center}
\small
\begin{tabularx}{\textwidth}{P{0.18\textwidth}P{0.28\textwidth}X}
\toprule
Paper quantity & \(\mathcal L\) counterpart & Quantitative relation \\
\midrule
\(\Time_D(n)\) & decider fuel \(f_D(n)\) & Term to TM:
\(\Time_D(n)\leq C_L(|D_{\mathcal L}|+f_D(n)+1)^6\).  TM to term:
fuel \(c_DT(T+|u|)\). \\
\(\Time_S(n)\) & maximum sampler-mode fuel \(f_S(n)\) & Same two bounds,
taking the maximum over the dimension, marginal, linear, and factor calls. \\
\(|D|\) & \(|\operatorname{ser}(D_{\mathcal L})|\) & Each translation is
linear: at most \(a_0|D|+b_0\) or \(a_U|D_{\mathcal L}|+b_U\). \\
\(|S|\) & \(|\operatorname{ser}(S_{\mathcal L})|\) & The same linear bounds;
static CL tables and source descriptions are counted once. \\
\(\lambda\)-bounded & size at most \(\lambda\), fuel at most \(n^\lambda\) &
Preserved in both directions as \(A\lambda\)-bounded for a universal integer
\(A\), as proved below. \\
Threshold \(n\geq C_0\) & same semantic index condition & Translation does not
change \(n\); resource estimates may replace it by
\(n\geq\max\{C_0,2\}\), never hide it in \(\poly\). \\
\bottomrule
\end{tabularx}
\end{center}

For the paper, the next theorem buys invariance of the verifier class under a
change of programming language, after one universal rescaling of \(\lambda\).
\begin{theorem}[Preservation of \(\lambda\)-boundedness]
\label{thm:lambda-preservation}
Fix the encodings and primitive registry above.  There is a universal integer
\(A\geq1\) such that:
\begin{enumerate}
\item translating any \(\lambda\)-bounded paper verifier term-by-term with
Theorem~\ref{thm:tm-to-l} gives an \(A\lambda\)-bounded
\(\mathcal L\)-verifier;
\item compiling any \(\lambda\)-bounded \(\mathcal L\)-verifier with
Theorem~\ref{thm:l-to-tm} gives an \(A\lambda\)-bounded paper verifier.
\end{enumerate}
The translation preserves the verifier's index and its input/output function.
\end{theorem}

\begin{proof}
For the first direction, \(|M|\leq\lambda\) and the size theorem gives
\(|\langle\!\langle M\rangle\!\rangle|\leq(a_0+b_0)\lambda\), since \(\lambda\geq1\).
At index \(n\geq2\), the paper machine runs for \(T\leq n^\lambda\).  It can
inspect or emit at most \(T\) tape symbols.  Across at most six input tapes the
portion relevant to the computation therefore has encoded length at most
\(8n^\lambda\), after increasing the constant to cover \(\log n\) and mode
tags.  Also
\(c_M\leq c_0(\lambda+7)\leq8c_0n^\lambda\).  Theorem~\ref{thm:tm-to-l}
then requires at most
\[
 8c_0n^\lambda\cdot n^\lambda\cdot9n^\lambda
 =72c_0n^{3\lambda}
 \leq n^{(3+\lceil\log_2(72c_0)\rceil)\lambda}.
\]
The last inequality uses \(n\geq2\) and \(\lambda\geq1\).  It applies to
both the sampler and decider.

For the second direction, hardwiring gives description size at most
\((a_U+b_U)\lambda\).  With fuel at most \(n^\lambda\), the lazy-input form of
\eqref{eq:decider-term-time} is at most
\[
 C_L(\lambda+n^\lambda+1)^6
 \leq C_L(3n^\lambda)^6
 \leq n^{(6+\lceil\log_2(C_L3^6)\rceil)\lambda}.
\]
We used \(\lambda\leq2^\lambda\leq n^\lambda\).  The same calculation holds
for every sampler mode.  Taking \(A\) to be the maximum of the four displayed
description and exponent coefficients proves both assertions.  Semantics is
preserved by Theorems~\ref{thm:tm-to-l} and~\ref{thm:l-to-tm}; neither
translation changes the binary index.
\end{proof}

The constants \(a_0,b_0,c_0,a_U,b_U,C_L\), fixed arities, serialization tags,
and the universal simulation exponent are absorbed when the paper writes
\(O(\cdot)\) or \(\poly(\cdot)\): they are universal after the two languages
are fixed.  The particular \(|D|,|S|\), a supplied \(\lambda\), a fuel or time
parameter, and hardwired \(|M|\) are \emph{arguments} of those polynomials and
cannot be hidden as universal constants.  Nor can the semantic threshold
\(C_0\), the change from \(\lambda\) to \(A\lambda\), or constants appearing
inside completeness and soundness inequalities be erased by \(\poly\)
notation.

\section{Self-reference as a construction on descriptions}
\label{sec:self-reference}

\begin{reminder}
Kleene's theorem in Section~\ref{sec:refresher-kleene} produced a behavioral
fixed point by applying a parameter-fixing compiler to its own code.  We now
repeat that construction with explicit description-size bounds and identify
it with the paper's diagonal verifier.
\end{reminder}

\subsection{The recursion theorem with size accounting}

Let \(\varphi_e^{(k)}\) be the partial function of the \(k\)-input paper
machine with description \(e\).  We use equality of partial functions only in
the conclusion; every operation in the construction acts on strings.

\begin{lemma}[The \(s\)-\(m\)-\(n\) operation]
\label{lem:smn}
For every \(r,k\geq1\) there is a total source transformer \(s_{r,k}\) such
that
\[
 \varphi_{s_{r,k}(e,z_1,\ldots,z_r)}^{(k)}(x)
 =\varphi_e^{(r+k)}(z_1,\ldots,z_r,x).
\]
Under the paper's hardwiring convention there are constants
\(a_s,b_s\), depending only on \(r,k\), for which
\[
 |s_{r,k}(e,z)|\leq a_s\bigl(|e|+\textstyle\sum_i|z_i|\bigr)+b_s,
\]
and the output code is computed in time linear in the right-hand side.
\end{lemma}

\begin{proof}
Output the fixed wrapper which writes
\(\operatorname{enc}_{r+k+1}(e,z_1,\ldots,z_r,x)\) on its work tape and invokes
the universal machine \(U_{r+k}\).  The wrapper contains one copy of each
static bit and constant code for tuple encoding and universal simulation.  It
therefore has the displayed linear length and can be printed in linear time.
The universal-machine theorem gives the stated partial-function equation on
halting inputs; on a nonhalting input both sides simulate the same machine
forever.
\end{proof}

For the paper, the next theorem buys self-reference while keeping the fixed
point's source length linear in the source transformer.
\begin{theorem}[Kleene's second recursion theorem, paper model]
\label{thm:kleene-two}
Let \(Q\) be a total computable transformation from descriptions of
\(k\)-input machines to descriptions of \(k\)-input machines.  There is an
effectively constructible description \(e_Q\) such that
\[
 \varphi_{e_Q}^{(k)}=\varphi_{Q(e_Q)}^{(k)}.
\tag{10.1}\label{eq:kleene-fixed}
\]
If \(q\) is a description of the machine computing \(Q\), then, for constants
\(a_K,b_K\) fixed by \(k\) and the encodings,
\[
 |e_Q|\leq a_K|q|+b_K,
\]
and \(e_Q\) is constructed in time polynomial in \(|q|\).
\end{theorem}

\begin{proof}
By Lemma~\ref{lem:smn}, let \(s(z,z)\) denote the description obtained by
hardwiring two copies of \(z\) into the first two inputs of a suitable
\((k+2)\)-input machine.  Define the partial computable function
\[
 G(z,x)=\varphi_{Q(s(z,z))}^{(k)}(x).
\]
A machine for \(G\) first computes the finite string \(s(z,z)\), runs the total
code transformer \(Q\) on it, and universally evaluates the resulting
\(k\)-input code on \(x\).  Let \(p\) be its description and set
\(e_Q=s(p,p)\).  Then for every \(x\), including the case in which the final
computation diverges,
\[
 \varphi_{e_Q}^{(k)}(x)
 =\varphi_p^{(k+1)}(p,x)
 =\varphi_{Q(s(p,p))}^{(k)}(x)
 =\varphi_{Q(e_Q)}^{(k)}(x),
\]
which proves \eqref{eq:kleene-fixed}.

The program \(p\) contains one copy of \(q\) and fixed code for \(s\) and the
universal evaluator, so \(|p|\leq a_1|q|+b_1\).  Applying the linear size bound
of Lemma~\ref{lem:smn} to \(s(p,p)\) gives
\(|e_Q|\leq2a_s a_1|q|+(2a_sb_1+b_s)\).  These are the asserted
\(a_K,b_K\).  Printing the two wrappers and executing the total hardwiring
transformations takes polynomial time under the paper's conventions.
\end{proof}

\subsection{The call-by-value Z combinator and YCode}

\begin{reminder}
Raw \(Y\) unfolds too early under call-by-value; Section~\ref{sec:refresher-lambda}
derived the eta-delayed \(Z\) combinator.  The constructor below realizes that
delay as a typed, constant-size self-code link.
\end{reminder}

For orientation, call-by-value lambda calculus uses
\[
 Z=\lambda f.(\lambda x.f(\lambda u.xx u))
                 (\lambda x.f(\lambda u.xx u)),
\]
rather than the call-by-name \(Y\), because eager evaluation of \(Yf\) unfolds
before supplying an argument.  This untyped term has constant size, but it
does not itself provide a canonical code value of the fixed point.  Our typed
description constructor does.

For a partial body \(P\) with its one distinguished quoted hole, define
\[
 \operatorname{YCode}(P):=\operatorname{Fix}(P),\qquad
 d_P:=\Quote(\operatorname{YCode}(P)).
\]
The source equation is
\[
 \operatorname{unfold}(\operatorname{YCode}(P))
 =\Specialize(P,\{\mathsf{self\_code}\mapsto d_P\}).
\tag{10.2}\label{eq:ycode-source}
\]
It is an equation between finite descriptions.  Its right side is produced on
demand by Lemma~\ref{lem:specialization}; the runtime evaluator instead uses a
constant-size environment link.

For the paper, the next theorem buys an operational self-code equation whose
unfolding cost and serialized size are both explicit.
\begin{theorem}[Description-level fixed-point equation]
\label{thm:ycode}
For every well-sorted \(P\), argument tuple \(u\), and \(f\geq c_Y=3\),
\[
 \operatorname{eval}_{\mathcal L}(\operatorname{YCode}(P),u;f)
 =\operatorname{eval}_{\mathcal L}\!\left(
   \Specialize(P,\{\mathsf{self\_code}\mapsto d_P\}),u;f-c_Y\right).
\tag{10.3}\label{eq:ycode-eval}
\]
Equivalently, if \(f_P\) is the code transformer represented by \(P\), the
right side is conventionally written
\(\operatorname{eval}_{\mathcal L}(f_P(\operatorname{YCode}(f_P)),u;
f-c_Y)\).  Moreover
\[
 |\operatorname{YCode}(P)|=|P|+c_{\rm fix}
\]
for the fixed four-bit Fix tag and its fixed sort annotation
\(c_{\rm fix}=O(1)\).
\end{theorem}

\begin{proof}
The unique Fix application rule performs three administrative actions: push
the application frame, install the self-code heap link, and enter \(P\).  The
resulting control, ordinary environment, and continuation are the same as for
the specialized body; a lookup of the sole hole follows the installed link and
therefore returns \(d_P\).  Subsequent transitions coincide by
Theorem~\ref{thm:l-determinism}.  Exactly three units have been consumed, which
proves \eqref{eq:ycode-eval}, including matching out-of-fuel and type-error
outcomes.  Canonical serialization of Fix writes its fixed tag and annotation
followed by one copy of the serialization of \(P\); it never expands the
self-code link.  This proves the size equality.
\end{proof}

\subsection{The paper's explicit diagonal program}

\begin{reminder}
The expression \(\mathcal F(\mathcal F,M,L,\ldots)\) is the recursion theorem
unrolled: hardwiring the first copy of \(\mathcal F\) manufactures the
five-input decider whose code is then supplied to \(\Compress\).
\end{reminder}

Let \(M:P\{\mathsf{MachineDesc}\}\) and
\(L:P\{\mathsf{Level}\}\) be closed literal terms, let \(S_L\) be the sampler
description returned by \(\ComputeSampler(L)\), and let \(n,x,y,a,b\) be the
five bound inputs.  The partial body corresponding to \code{fig:halt_f} is
\[
\begin{split}
 \Psi_{M,L}=\lambda(n,x,y,a,b).\ \operatorname{If}(&
  \operatorname{Prim}(\mathsf{halts\_within},M,n),\ \mathsf{true},\\
 &\operatorname{Eval}(
   \operatorname{ans}(\operatorname{Eval}(
      \Quote(\Compress),\\[-2pt]
 &\quad (\operatorname{quoted\_pair}(\Quote(S_L),
       \operatorname{Hole}(\mathsf{self\_code},
          \mathsf{Quoted}(\mathsf{Decider}))),L),F_C)),\\[-2pt]
 &\quad(n,x,y,a,b),\operatorname{FuelBound}(n,L))).
\end{split}
\tag{10.4}\label{eq:psi-ml}
\]
The finite-fuel symbol \(F_C\) is any proved bound for the terminating
\(\Compress\) call; \(\operatorname{ans}\) selects the returned decider
description.  Define
\[
 D_{M,L}=\operatorname{YCode}(\Psi_{M,L}),\qquad
 V_{M,L}=(S_L,\Quote(D_{M,L})).
\tag{10.5}\label{eq:dml}
\]

For the paper, the next theorem buys the identification of the published
\(\mathcal F\) trick, Kleene's theorem, and the IR fixed-point constructor.
\begin{theorem}[Equivalence of the three fixed-point presentations]
\label{thm:fixed-point-equivalence}
The following are translations of the same description-level construction:
\begin{enumerate}
\item the paper's explicit
\(\mathcal F(\overline{\mathcal F},\overline M,L,n,x,y,a,b)\);
\item the Kleene fixed point of the transformer which hardwires
\((\overline M,L)\) and inserts the resulting decider code into \(\Compress\);
\item the closed description \(D_{M,L}\) in \eqref{eq:dml}.
\end{enumerate}
Under the simulations of Section~\ref{sec:simulation}, their deciders compute
the same partial function.  Their description lengths are all
\(O(|M|+\log L+1)\), and their running times differ by a polynomial in
\(|M|,\log L,n\), the invoked fuel bounds, and the invoked machine times.
\end{theorem}

\begin{proof}
In the paper, Step 2 of \(\mathcal F\) applies the hardwiring operation to the
eight-input code supplied in its first argument.  When that argument is
\(\overline{\mathcal F}\), the resulting five-input code is precisely the code
of the outer decider
\(D(n,x,y,a,b)=\mathcal F(\overline{\mathcal F},\overline M,L,n,x,y,a,b)\).
Thus the string given to \(\Compress\) is its own decider description.  This
is the concrete \(s\)-\(m\)-\(n\) operation of Lemma~\ref{lem:smn} with the
diagonal argument already chosen.

Alternatively, let \(Q(e)\) output the code which performs the halting test and,
on failure, applies \(\Compress\) to \((S_L,e)\) and runs the returned decider.
Theorem~\ref{thm:kleene-two} gives \(e_Q\) with
\(\varphi_{e_Q}=\varphi_{Q(e_Q)}\), which is exactly the behavior just
described.  Equation~\eqref{eq:ycode-source} performs the same substitution
with \(e_Q\) represented by \(d_P\); Theorem~\ref{thm:ycode} proves the
operational equation for \(D_{M,L}\).  Translating it to a machine with
Theorem~\ref{thm:l-to-tm}, or translating \(\mathcal F\) with
Theorem~\ref{thm:tm-to-l}, preserves that behavior.

The displayed term contains one copy of \(M\), one binary literal \(L\), and
fixed code for the remaining operations.  Theorem~\ref{thm:ycode} therefore
gives \(|D_{M,L}|=O(|M|+\log L+1)\).  The paper's hardwiring and the size
calculation in Theorem~\ref{thm:kleene-two} give the same bound for the first
two forms.  Finally, the two simulation theorems replace each bounded call by a
fixed polynomial overhead.  Composition of finitely many such polynomials is
a polynomial in the listed quantities, proving the time statement.
\end{proof}

The paper notes that an earlier version used Kleene's recursion theorem and
that its published presentation uses \(\mathcal F\) explicitly.
\gt{fig:halt_f}{ground-truth/gt-12-compression.tex}  Theorem~\ref{thm:fixed-point-equivalence}
explains why those presentations differ in notation rather than in the
description passed to compression.

What must \(\Compress\) do with the quoted body?  It first applies
\(\Specialize\) to the static verifier and parameter holes, obtaining closed
syntax with the size bound of Lemma~\ref{lem:specialization}.  It then traverses
that syntax to build the descriptions for introspection, answer reduction, and
repetition, and it records the resulting byte lengths and fuel bounds.  It
never asks whether two terms compute the same function.  Such extensional
equality is neither available nor needed.  This self-reference is not the hard
part: it supplies no low-degree soundness, rigidity, repetition, or gap
preservation.  Those remain the substantive cited theorems.

\section{Cook--Levin for \texorpdfstring{\(\mathcal L\)}{L}-descriptions}
\label{sec:l-cook-levin}

\begin{reminder}
A machine run is a path of configurations under one local transition rule;
Section~\ref{sec:refresher-tm} introduced this viewpoint.  Here the same idea
is applied to the fixed universal evaluator, turning a fuel-bounded lambda run
into locally checkable tableau data.
\end{reminder}

\subsection{A local encoding of bounded evaluation}

For closed \(t\), input \(u\), and fuel \(F\), the \emph{bounded trace} is
\[
 \mathcal T(t,u,F)=(C_0,C_1,\ldots,C_r),\qquad r\leq F,
\]
where \(C_0\) is the initial CEK configuration and each \(C_{i+1}\) is its
unique small-step successor.  A row records the current term pointer,
environment-frame pointers, continuation, heap, primitive microstate, outcome
flag, and binary fuel.  If evaluation halts early, a fixed self-looping final
state pads the trace to \(F+1\) rows.  This is the promised sequence of
\((\text{term},\text{environment},\text{fuel})\) configurations, with the
continuation and heap made explicit so that one-step validity is local.

To obtain a finite-alphabet tableau, compile the CEK evaluator and each charged
primitive to a fixed multitape Turing machine \(\widehat U_{\mathcal L}\).
The read-only code track contains \(\operatorname{ser}(t)\); input tracks hold
lazy argument tapes; work tracks hold the environment, heap, stack, and fuel.
Pointer arithmetic, traversal, and a primitive of charge \(k\) are expanded
into their elementary head moves.  One microstep changes the finite control,
the cells under the heads, and the head positions by at most one.

\begin{lemma}[Transition-window lemma]
\label{lem:transition-window}
There is a finite alphabet \(\Gamma_{\mathcal L}\), a fixed number \(r_0\) of
tracks, and a Boolean function
\[
 \tau:\bigl(\Gamma_{\mathcal L}^{3}\bigr)^{r_0}
       \times Q_{\mathcal L}\longrightarrow
       \Gamma_{\mathcal L}^{r_0}\times Q_{\mathcal L}
\]
with the following property.  For every nonboundary cell \(j\), the symbols in
cell \(j\) of row \(i+1\), together with the next finite-control state, are
determined by the radius-one windows
\((j-1,j,j+1)\) of row \(i\) and the current control state.  Boundary cells use
a fixed endmarker version of the same rule.  Conversely, if two consecutive
rows have one head marker on each track, agree with \(\tau\) at every window,
and have valid endmarkers, then the second row is exactly the successor of the
first under the small-step semantics.
\end{lemma}

\begin{proof}
Encode a head by pairing the tape symbol with one of the finitely many machine
states when that head is the active head, and use a separate finite-control
track for the other heads.  In one transition, a cell farther than one position
from every head is copied.  At a head cell, the transition table determines the
written symbol and next state from the old state and scanned symbols.  At its
left and right neighbors, the table determines whether the head marker enters;
it cannot enter any other cell because heads move by at most one.  These cases
depend only on the three displayed cells on each track.  They are mutually
exclusive once the old row has one head per track, and therefore define the
finite function \(\tau\).  Endmarkers replace missing neighbors by a fixed
boundary symbol.

For the converse, consider each track.  Away from the unique old head, the
local copy case forces equality.  In the three-cell neighborhood of that head,
the transition case forces the unique symbol write, direction, and new head
position specified by the machine table.  The local rules also force all other
cells not to acquire a head.  Repeating this argument for every track and
using the shared finite-control component gives precisely one transition of
\(\widehat U_{\mathcal L}\).  Primitive execution and CEK operations introduce
no exceptional case because they were compiled into the same elementary
machine steps.  The correspondence between those microstates and the weighted
semantics was part of the construction in Section~\ref{sec:resource-lambda}.
\end{proof}

\subsection{The succinct 3SAT circuit}

Let the row width be \(W\).  In \(F\) microsteps no head visits beyond the
initial code and input plus \(F\) cells, so
\(W=O(|t|+|u|+F)\).  Introduce Boolean variables for the constant-size binary
encoding of every tableau cell and control-state bit.  The formula asserts:
the initial row contains \(t,u,F\); every row and track has the prescribed
format and head marker; each adjacent pair obeys every transition window; and
the last row is accepting.  Each assertion involves a constant window except
binary addressing and the initial read-only strings.  Fixed Tseitin gadgets
turn it into 3CNF with \(O(FW)\) variables and clauses.

For the paper, the next theorem buys a small circuit describing an enormous
bounded verifier tableau, which is the input expected by answer reduction.
\begin{theorem}[Succinct Cook--Levin for \(\mathcal L\)]
\label{thm:l-succinct-sat}
There is an algorithm which, on
\((t,n,F,Q,\sigma,x,y)\), where \(|t|\leq\sigma\) and \(|x|,|y|\leq Q\),
outputs a Boolean circuit \(C_{t,n,F,Q,\sigma,x,y}\) deciding membership of an
indexed signed clause in a 3CNF formula \(\Phi\).  Its index width is
\(O(\log F+\log\sigma)\), its size and construction time are
\[
 \poly(\log n,\log F,Q,\sigma),
\tag{11.1}\label{eq:l-succinct-size}
\]
and for every pair of answer strings \(a,b\) the formula is satisfiable with
the prescribed input bits if and only if
\(\Eval_{\mathcal L}(\Quote(t),(n,x,y,a,b);F)\) accepts.  Holding the public
input fixed, the transition component of the circuit has size
\(\poly(\log F,|t|)\); the \(Q\) term in \eqref{eq:l-succinct-size} is solely
the hardwired initialization data \(x,y\).
\end{theorem}

\begin{proof}
Given a clause index, the circuit first decodes whether the requested clause is
an initialization, uniqueness/format, transition, or acceptance clause.  Row
and column numbers use \(O(\log F+\log W)\) bits.  For a transition clause it
computes the addresses of a radius-one window and evaluates the finite truth
table \(\tau\) from Lemma~\ref{lem:transition-window}; binary addition,
comparison, and multiplexing have size polynomial in those address lengths.
For an initialization clause it selects a bit of the canonical code or public
input.  A balanced selector for a hardwired string has size linear in
\(|t|+Q+\log n\).  Format and final-state clauses use fixed truth tables.
Finally it emits the three variable addresses and signs of the selected
3CNF gadget.  This proves the circuit and construction bounds.  Since
\(W=O(|t|+Q+F)\), \(\log W=O(\log F+\log\sigma+\log(Q+1))\), which is absorbed
by \eqref{eq:l-succinct-size} under \(Q\leq F\), as in the paper.

If evaluation accepts, assign every tableau variable from its unique bounded
trace and every Tseitin variable from its gadget value.  Initialization,
transition, and acceptance clauses are then satisfied.  Conversely, any
satisfying assignment contains the prescribed initial row.  Induction on the
row number and the converse part of Lemma~\ref{lem:transition-window} forces
each later row to be the unique evaluator successor.  The final clauses force
that trace to accept.  The prescribed answer bits occupy the initial answer
tapes, so the equivalence holds for each \(a,b\).
\end{proof}

This is the \(\mathcal L\)-description restatement of
\code{prop:standard-succinct-sat}.  The paper's version takes \(D,n,T,Q,
\sigma,x,y\), reserves the first \(2T\) encoded positions for each answer, and
has size \(\poly(\log n,\log T,Q,\sigma)\).
\gt{prop:standard-succinct-sat}{ground-truth/gt-10-answer-reduction.tex}
Theorem~\ref{thm:l-to-tm} gives the same proposition by compilation; the proof
above shows directly why a lambda evaluation trace has the required locality.

More explicitly, use the paper's tape-alphabet encoding
\(\operatorname{enc}_\Gamma(0)=00\),
\(\operatorname{enc}_\Gamma(1)=01\),
\(\operatorname{enc}_\Gamma(\sqcup)=10\), and reserve \(2F\) bits in the
tableau assignment for each answer tape.  For every
\(a,b\in\{0,1\}^{2F}\), the resulting formula has an extension \(c\) if and
only if there are prefixes \(a',b'\), of lengths at most \(F\), such that
\[
 a=\operatorname{enc}_\Gamma(a',\sqcup^{F-|a'|}),\qquad
 b=\operatorname{enc}_\Gamma(b',\sqcup^{F-|b'|}),
\]
and the term accepts \((n,x,y,a',b')\) within fuel \(F\).  The forward
direction follows because the initialization clauses force the two reserved
blocks to be valid padded tape rows; the backward direction fills those blocks
from the accepting trace.  This is the exact placement property needed when
the two blocks are separated in \eqref{eq:decoupled-five}.

\subsection{From shared 3SAT data to decoupled 5SAT}

The paper removes two undesirable features before low-degree encoding.  A
3SAT clause reads three locations of one tableau assignment, while three
low-degree answers need not be restrictions of one assignment.  Also, the
strings \(a,b\) are embedded inside that assignment, whereas the later sampler
must query them as independent blocks.  Make five assignments
\(a,b,u_3,u_4,u_5\) and impose clauses of the form
\[
 a_{i_1}^{o_1}\vee b_{i_2}^{o_2}\vee
 (u_3)_{i_3}^{o_3}\vee (u_4)_{i_4}^{o_4}\vee
 (u_5)_{i_5}^{o_5}.
\tag{11.2}\label{eq:decoupled-five}
\]
The original 3SAT clauses are applied to \(u_3\); constant-size five-literal
gadgets enforce \(u_3=u_4=u_5\), copy the first \(2F\) encoded answer bits to
\(a,b\), and pad unused positions by the alternating blank encoding.  Each
gadget is indexed by comparisons and additions on
\(O(\log F+\log\sigma)\)-bit addresses, so it adds only that many gates.
Therefore the succinct decoupled circuit retains the bound
\eqref{eq:l-succinct-size}.  Its padded form makes all five blocks have one
common power-of-two length.  A satisfying five-tuple restricts to a satisfying
trace assignment, and a satisfying trace extends by copying and padding; this
proves the required if-and-only-if.
\gt{sec:succinct-deciders}{ground-truth/gt-10-answer-reduction.tex}

The intended analytic chain is
\[
\begin{split}
\mathtt{Quoted\{Decider\}}
&\xrightarrow{\ \mathtt{bounded\_trace}\ }
 \mathtt{BoundedTrace}\\
\mathtt{BoundedTrace}
&\xrightarrow{\ \mathtt{cook\_levin}\ }
 \mathtt{Succinct3SAT}\\
\mathtt{Succinct3SAT}
&\xrightarrow{\ \mathtt{decouple5}\ }
 \mathtt{SuccinctDecoupled5SAT}
\\
\mathtt{SuccinctDecoupled5SAT}
&\xrightarrow{\ \mathrm{padding}\ }
 \mathtt{PaddedSuccinctDecoupled5SAT}.
\end{split}
\]
These are the transformations planned for TB3 in \code{briefs/23-tb3.md}.
That brief is planned work, not evidence that the functions exist or that any
claim has been checked.  At the time of this note, no \code{bounded_trace},
\code{cook_levin}, or \code{decouple5} implementation is present in
\code{src/}; the general contracts remain \grade{CITED}.

\subsection{Tseitin and arithmetization}

\begin{reminder}
At this point the quoted verifier has already been reduced to a finite circuit
describing local computation checks.  The transformations below manipulate
that circuit description; they are computable source maps, while their
low-degree and quantum soundness consequences are separate analytic claims.
\end{reminder}

Given a Boolean circuit \(C\) with gates \(g_i\), Tseitin introduces a wire
variable \(w_i\) and an equivalence formula
\[
 z_i=(g_i(x,w)\wedge w_i)\vee
     (\neg g_i(x,w)\wedge\neg w_i),
 \qquad F=\bigwedge_i z_i.
\]
The required contract is
\(C(x)=1\) iff there exists \(w\) with \(F(x,w)=1\), with size and
construction time \(O(|C|)\).\gt{def:tseitin}{ground-truth/gt-10-answer-reduction.tex}
The project follows the explicit gate-equivalence form in
\code{ground-truth/nw19/nw19-tseitin-arith.tex} and also conjoins the output
literal \(w_{\rm out}\), so the stored formula explicitly asserts acceptance.
That convention is marked \grade{SOURCE\_REPAIR}; it is not silently attributed
to the source.

The Julia constructors \code{Circuit}, \code{Input}, \code{Gate},
\code{NotGate}, \code{AndGate}, and \code{OrGate} make the finite DAG and its
topological order explicit.  The function \code{tseitin} returns
\code{Checked{TseitinFormula,CertNode}}.  Its replay recomputes the formula's
occurrence vector; it does not, by that replay alone, establish the general
Cook--Levin or Tseitin semantic equivalence.  Those implications are exhausted
only on the small TB0 fixture and remain cited in general.

Arithmetization recursively applies
\[
  \neg A\mapsto 1-A,\qquad
  A\wedge B\mapsto AB,\qquad
  A\vee B\mapsto A+B-AB.
\]
It agrees with the Boolean formula on the Boolean cube.
\gt{def:formula-arithmetization}{ground-truth/gt-10-answer-reduction.tex}
The related NW19 rule is \code{def:arithmetization} in
\code{ground-truth/nw19/nw19-tseitin-arith.tex}.  Julia's \code{arith_q}
returns a sparse formal \code{Poly}; its \code{ArithTseitin} certificate checks
support degrees against a structural bound.

\subsection{The occurrence-vector discrepancy}

There is a discrepancy in the individual-degree bookkeeping as we read NW19;
our reading may be wrong.  The source proposition states individual degree at
most two for the Tseitin arithmetization.
\gt{prop:tseitin-arith-degree}{ground-truth/gt-10-answer-reduction.tex}
Reading the NW19 formula literally as a formula tree, arithmetization is
multilinear in leaf occurrences, not necessarily in circuit wires.  The same
wire can occur in several gate gadgets.  With the project's explicit output
literal, the count we compute is
\[
 \operatorname{occ}(x_j)=2\operatorname{fanout}(x_j),\qquad
 \operatorname{occ}(w_i)=2+2\operatorname{fanout}(w_i)
       +\mathbf 1_{i=\mathrm{out}},
\]
and structural induction gives
\[
 \deg_v(F_{\rm arith})\leq \operatorname{occ}_v(F).
\]
The function \code{tseitin_occurrence_account} computes this vector;
\code{occurrences} independently traverses the formula; and \code{arith_q}
compares the structural account with actual sparse support.

Claim C8 is only \claim{C8}{TESTED}: the two-gate regression gives
\((2,2,2,4,3)\), and the 16-variable TB0 formula attains
\[
(2,0,0,0,2,2,0,0,0,0,6,4,4,4,4,3).
\]
The general occurrence formula is a written derivation, not a machine theorem.
The downstream soundness repair is only \claim{C5}{SKETCH}.  With a symbolic
bound \(\delta_F=\max_v\operatorname{occ}_v(F)\), the conservative formula-test
term becomes \((\delta_F+5d)m'/q\), and
\code{P_formula_structural} is an additional obligation.  Under an explicit
fanout-reduction hypothesis, \(\operatorname{fanout}_{\max}\leq2\) gives
\(\delta_F\leq7\) with the output literal, and \(d\geq8\) suffices for the
individual-degree construction bound.  Absorption into the asymptotic
parameter choice still needs the stated growth hypotheses; it is not promoted
by the two finite computations.

\section{The low-degree PCP as algebra}

\subsection{Interpolation and the vanishing polynomial}

For \(M=2^m\) and \(a=(a_y)_{y\in\{0,1\}^m}\), define
\[
 \operatorname{ind}_{m,y}(x)=
 \prod_{j:y_j=1}x_j\prod_{j:y_j=0}(1-x_j),\qquad
 g_a(x)=\sum_y a_y\operatorname{ind}_{m,y}(x).
\]
Then \(g_a(y)=a_y\) on the Boolean cube and \(a\mapsto g_a\) is linear.
This is the low-degree code: the assignment is represented by the evaluation
table of its unique multilinear extension.
\gt{sec:ld-encoding}{ground-truth/gt-03-prelim.tex}
The Julia function \code{g_a} constructs the sparse polynomial; block
coordinates are carried by \code{VarLayout} and \code{VarBlock};
\code{dependency_blocks} recomputes the actual support dependency.

For five satisfying blocks, let \(g_i\) be their multilinear extensions and
write \(z=(x_1,\ldots,x_5,o,w)\in\F_q^{m'}\).  The central analytic object is
\[
 c_0(z)=F_{\rm arith}(z)
        \prod_{i=1}^{5}\bigl(g_i(x_i)-o_i\bigr).
\]
On a Boolean point, either \(F_{\rm arith}=0\), or the indicated clause is
present and at least one factor \(g_i(x_i)-o_i\) vanishes.  Hence \(c_0\)
vanishes on \(\{0,1\}^{m'}\).  The code executes the product in
\code{build_c0}; the \code{BuildC0} certificate replays the structural and
actual degree accounts.  The implication from a satisfying decoupled instance
to Boolean-cube vanishing is checked only on the TB0 fixture and cited in
general.\gt{thm:pcp-decider}{ground-truth/gt-10-answer-reduction.tex}

\subsection{Division by the Boolean ideal}

Let \(z_i(1-z_i)=z_i-z_i^2\).  The ground-truth basis proposition says that a
degree-bounded polynomial vanishing on the Boolean cube lies in the ideal
generated by these \(m'\) polynomials.
\gt{prop:zero-basis}{ground-truth/gt-10-answer-reduction.tex}
The implementation uses the following deterministic monomial rewrite.

\begin{lemma}[Executable zero-basis rewrite]
For a coefficient \(a\), a monomial \(u\) independent of the displayed power,
and \(e\geq2\),
\[
 a u z_i^e
 \longmapsto
 a u z_i^{e-1}-a u z_i^{e-2}\,z_i(1-z_i).
\]
Iterating in a fixed variable and monomial order yields polynomials \(c_i\) and
a remainder \(r\), multilinear in every variable, such that
\[
 c_0=\sum_{i=1}^{m'}c_i z_i(1-z_i)+r.
\]
If the coordinate degree vector of \(c_0\) is \(D\), then
\(\deg_j(c_i)\leq D_j\), \(\deg_i(c_i)\leq\max(D_i-2,0)\), and after reducing
coordinates \(1,\ldots,i-1\), the running remainder has degree at most one in
those coordinates.
\end{lemma}

\begin{proof}
Expanding the right-hand side of one rewrite gives
\[
 a u z_i^{e-1}-a u z_i^{e-2}(z_i-z_i^2)=a u z_i^e.
\]
Thus every step preserves the coefficient polynomial after adding the recorded
quotient contribution.  It lowers the current exponent to one and places
power \(e-2\) in the quotient, proving the degree bounds.  Induction over the
ordered coordinates gives the displayed decomposition and a final
multilinear remainder.  If \(c_0\) vanishes on the Boolean cube, then so does
\(r\); uniqueness of multilinear interpolation on \(\{0,1\}^{m'}\) makes
\(r=0\).  This is the constructive division used in the proof of
\code{prop:zero-basis} in
\code{ground-truth/gt-10-answer-reduction.tex}.
\end{proof}

\code{zero_basis_decompose} stores every step as \code{RewriteStep} inside a
\code{ZeroDecomposition}.  \code{verify_rewrite_step} replays the scalar
identity, and \code{verify_zero_decomposition} reconstructs the entire right
side as a normalized coefficient dictionary.  Promotion to a zero certificate
also requires empty remainder.  This is stronger than checking the equality at
sampled points, but it remains an instance certificate, not a proof of the PCP
soundness theorem.

\subsection{PCP view and local verifier}

A \emph{PCP view} at \(z\) is
\[
 \operatorname{ev}_z(\Pi)=
 (g_1(x_1),\ldots,g_5(x_5),c_0(z),c_1(z),\ldots,c_{m'}(z)).
\]
The proof object is the tuple of evaluation tables for those polynomials.
The labels are \code{def:pcp-proof} and \code{def:pcp-eval} in
\code{ground-truth/gt-10-answer-reduction.tex}.  Julia represents them by
\code{PCPProof} and \code{PCPView}; \code{ev_z} refuses a \(g_i\) whose actual
support escapes block \(X_i\).

The verifier first reconstructs the succinct circuit, Tseitin formula, and
arithmetization.  It then checks exactly
\[
 \beta_0=F_{\rm arith}(z)\prod_{i=1}^5(\alpha_i-o_i),
 \qquad
 \beta_0=\sum_{j=1}^{m'}\beta_j z_j(1-z_j).
\]
These are the formula test and zero-on-subcube test.
\gt{fig:pcpverifier}{ground-truth/gt-10-answer-reduction.tex}
The Julia function \code{pcpverifier} executes these two equations on an
already supplied \code{TseitinFormula} and view.  In TB2 it does not implement
the figure's preceding reconstruction steps; the formula is construction-time
data.  \code{build_pcp} combines upstream certificates, coefficient division,
degree checking, and stored certified views.

\subsection{The four soundness layers}

The four layers must not be collapsed.

\paragraph{1. Algebraic soundness under a low-degree premise.}
Assume every submitted \(g_i,c_j\) has individual degree at most \(d\).  If the
formula-test polynomials are unequal, their difference has total degree at most
\((\delta_F+5d)m'\) under the occurrence-vector account.  If the zero-test
polynomials are unequal, their difference has total degree at most
\((2+d)m'\).  Schwartz--Zippel states that unequal polynomials of total degree
at most \(\Delta\) agree at a uniform point of \(\F_q^{m'}\) with probability
at most \(\Delta/q\).
\gt{lem:schwartz-zippel}{ground-truth/gt-03-prelim.tex}
Consequently
\[
 \Pr[\text{false formula identity}]
 \leq \frac{(\delta_F+5d)m'}{q},\qquad
 \Pr[\text{false zero identity}]
 \leq \frac{(2+d)m'}{q}.
\]
The paper's displayed first numerator uses \(2+5d\); the structural project
obligation substitutes the measured \(\delta_F\) under the discrepancy stated
above.  The field size enters nowhere mysteriously: \(q=2^k\) is chosen so
both fractions are below the acceptance threshold \(1/2\), together with the
other PCP parameter obligations.  The first structural inequality is extra to
the literal \code{def:pcpparams} predicate in the present project.

If acceptance is greater than \(1/2\) and both upper bounds are below
\(1/2\), each sampled equality is therefore a polynomial identity.  The zero
identity makes \(c_0\) vanish on the Boolean cube; the formula identity then
decodes five Boolean assignments satisfying the decoupled formula, and hence an
accepting bounded computation.  This conditional theorem is
\code{thm:pcp-decider} in
\code{ground-truth/gt-10-answer-reduction.tex}.  Its current project status is
\claim{C5}{SKETCH}, not CHECKED.

\paragraph{2. Enforcement of the low-degree premise.}
The simultaneous low-degree game queries points, axis-parallel lines, and
diagonal lines.  Its classical decider checks consistency of a point value with
a claimed line restriction.\gt{fig:ld-decider}{ground-truth/gt-07-ldt.tex}
The theorem extracting approximately consistent low-individual-degree
polynomial measurements from a successful quantum strategy is
\code{lem:ld-soundness} in \code{ground-truth/gt-07-ldt.tex}.  It is
\grade{CITED}.  Running \code{ld_decider} on honest finite inputs does not prove
this extraction theorem.

\paragraph{3. Polynomial identity from random evaluation.}
This is exactly the Schwartz--Zippel step above, used twice with two different
degree bounds.  The local program checks evaluations.  The analytic theorem
turns an acceptance probability into formal identity, provided the degree and
field-size premises hold.  No sample count can replace that implication.

\paragraph{4. Quantum consistency and rigidity.}
Answer reduction must connect different players' point and line answers to the
same polynomial measurements, control disturbance, and transport the result
back to the original verifier.  Those operator estimates use low-degree
soundness, typed consistency transformations, and later rigidity arguments.  They
are imported in the soundness contract of \code{thm:ar} in
\code{ground-truth/gt-10-answer-reduction.tex}; none is a sparse-polynomial
identity.

The computer adds exactly this: on a named instance it constructs the
polynomials, executes the division, compares normalized coefficients, computes
degree and dependency vectors, and evaluates the two local equations.  It does
not test the theorem.  Current statuses reflect that separation: C3 is
\textsc{TESTED}, C5 is \textsc{SKETCH}, C1 and C2 remain
\textsc{CONJECTURE}, and C8 is \textsc{TESTED} only on its two named circuits.

\section{The typed layer: conditionally linear functions}

\subsection{The inductive object and its laws}

\begin{definition}[Conditionally linear function]
Let \(V=\F^n\) with fixed coordinate-register subspaces.  The unique level-zero
conditionally linear (CL) function is zero.  A level-\(\ell\) CL function
chooses complementary registers \(V_1\oplus V_{>1}=V\), applies a linear map
\(A:V_1\to V_1\), and, conditional only on \(A(x^{V_1})\), selects a
level-\((\ell-1)\) CL function on \(V_{>1}\).  Its output is the sum of the
first-stage and continuation outputs.
\end{definition}
This is \code{def:cl-func} in
\code{ground-truth/gt-04-cl.tex}.  Julia mirrors it with
\code{AbstractCL}, \code{CLZero}, and \code{CLStep}.  The constructor verifies
disjoint factor/rest coordinate sets, a matrix acting inside the factor, a
fixed child level, and exact coverage of the rest register.  Thus an inhabitant
of \code{CLStep} has a constructed level; linearity is carried by the matrix,
not inferred from samples.

The marginal \(L_{\leq k}\) is the sum of the first \(k\) stage outputs,
together with their factor spaces and linear maps.  This is the decomposition
of \code{lem:cl-kth} in
\code{ground-truth/gt-04-cl.tex}.  Julia's \code{marginal_k} returns a
\code{CLMarginal}; \code{sum_stage_outputs} recomputes its value.

\begin{lemma}[Level laws]
Let \(L\) be level \(k\), let every continuation \(R_u\) be level \(\ell\),
and let \(L_j\) act on pairwise disjoint registers at levels \(\ell_j\).
Then
\[
 \operatorname{level}(\operatorname{concatenate}(L,R))=k+\ell,
 \qquad
 \operatorname{level}\!\left(\bigoplus_jL_j\right)=\max_j\ell_j.
\]
For CL distributions \(S_1,S_2\), their product, formed by direct-summing the
left maps and the right maps on independent seed registers, has level the
maximum of the two levels.
\end{lemma}

\begin{proof}
For concatenation, induct on the outer constructor of \(L\).  At level zero,
grafting \(R_0\) supplies exactly \(\ell\) stages.  At a \code{CLStep}, retain
its first linear stage and apply the induction hypothesis to each child; this
gives \(1+(k-1+\ell)=k+\ell\).  This is the construction in
\code{lem:cl-concat} of \code{ground-truth/gt-04-cl.tex}.

For direct sum, pad each component by zero stages to
\(r=\max_j\ell_j\).  At stage \(i\), use the direct sum of the components'
stage matrices on the direct sum of their factor registers.  The continuation
depends only on the tuple of earlier outputs, hence remains CL.  Induction on
\(r\) proves the law, matching \code{lem:cl-func-prod} in
\code{ground-truth/gt-04-cl.tex}.  A distribution product applies this law
separately to its left and right maps, so its level is the same maximum.
\end{proof}

The code functions are \code{concatenate}, \code{direct_sum},
\code{distribution}, and \code{product}.  Their type constructions carry the
level laws.  The stronger structural hypothesis that all compiler operations
of interest close inside one useful algebra is not implied by these individual
lemmas.

\subsection{Point and line maps}

For \(V=V_{\rm pt}\oplus V_{\rm coord}\oplus V_{\rm dir}\), write a seed as
\((u,s,v)\in\F_q^m\times\F_q\times\F_q^m\).  Let
\(L_v^{\rm lnf}\) be the canonical projection whose kernel is
\(\operatorname{span}(v)\), so \(L_v^{\rm lnf}u\) is a canonical representative
of the affine line \(u+\F_qv\).
\gt{def:cl-canonical}{ground-truth/gt-03-prelim.tex}
The line-representative use is \code{def:line-representative} in
\code{ground-truth/gt-07-ldt.tex}.  Define \(\chi(s)\) by splitting the fixed
integer representatives of \(\F_q\) into \(m\) equal buckets; this requires
\(m\mid q\).  With \(\pi_{i-1}\) zeroing the first \(i-1\) coordinates,
\[
\begin{aligned}
L_{\rm Point}(u,s,v)&=(u,0,0),\\
L_{\rm ALine}(u,s,v)&=(L_{e_{\chi(s)}}^{\rm lnf}u,s,0),\\
L_{\rm DLine}(u,s,v)&=(L_{\pi_{\chi(s)-1}(v)}^{\rm lnf}u,
                         s,\pi_{\chi(s)-1}(v)).
\end{aligned}
\]
Their levels are respectively one, two, and three by their stage
factorizations.\gt{sec:ld-game}{ground-truth/gt-07-ldt.tex}
The implementation names are \code{L_Point}, \code{L_ALine},
\code{L_DLine}, \code{L_lnf}, \code{chi}, and \code{pi_prefix}.  Since the
source's kernel-basis definition does not cover \(v=0\), the code uses
\(L_0^{\rm lnf}=I\) and marks the convention \grade{SOURCE\_REPAIR}.

The distribution \(\mu_{L,R}\) pushes one uniform shared seed through the two
maps.\gt{def:cl-dist}{ground-truth/gt-04-cl.tex}
For \((q,m)=(8,2)\), TB1 enumerates all \(8^5=32768\) seeds.  The
axis-line/point histogram has support 512, and the diagonal-line/point
histogram has support 18432, including the zero directions.  Both agree
entry-for-entry with separately transcribed formulas for
\code{lem:alnf} and \code{lem:dlnf} in
\code{ground-truth/gt-07-ldt.tex}; every marginal is replayed.  This finite
statement is \claim{C4a}{TESTED}.  It is not a statistical argument and does
not establish quantum low-degree soundness.

\subsection{Types and the six-copy sampler}

A \emph{typed sampler} is a finite type-graph-indexed family of CL maps sharing
one ambient seed space.  It first samples an oriented type edge and then pushes
one uniform seed through the selected left/right maps.
\gt{def:typed-sampler}{ground-truth/gt-06-types.tex}
Its induced distribution is \code{def:typed-sampler-sample} in the same file.
Julia's \code{TypedSampler} checks exact type coverage, a common field and seed
dimension, graph endpoints, and a padded common level.

\emph{Detyping} encodes the external graph choice into an untyped CL
distribution by a rejection-sampling construction.  It adds two levels,
preserves value-one completeness, and changes soundness error by
\(16^{|\mathrm{Type}|}\).\gt{lem:detyping-verifiers}{ground-truth/gt-06-types.tex}
The Julia function \code{detype} returns \code{CitedDetypedVerifier}; the
certificate is deliberately \grade{CITED}.  For the 54 answer-reduction types,
the displayed factor is \(16^{54}\), not an executed estimate.

The answer-reduction PCP sampler has the 18 types
\[
 \{\mathrm{Point}_i,\mathrm{ALine}_i,\mathrm{DLine}_i:1\leq i\leq6\}
\]
and a complete type graph.  Copies 1--5 query one \(m\)-variate \(g_i\) each;
copy 6 queries the bundle \((g_1,\ldots,g_5,c_0,\ldots,c_{m'})\) in dimension
\(m'\).  The six-copy construction and its ambient direct sums are in
\code{sec:ld-compiler} of
\code{ground-truth/gt-10-answer-reduction.tex}.  The current code uses
\code{PCPType}, \code{PCPRegisterLayout}, \code{PCPCLMap},
\code{PCPPaddedMap}, and \code{pcp_sampler}.  It carries a visible
\grade{SOURCE\_REPAIR} for the copy-6 coordinate-register conflict between the
displayed ambient space and the parser table.

\emph{Oracularization} is the typed transformation in which the isolated
roles receive the original left or right CL image while an oracle role receives
the common seed and can reconstruct both; the code name is
\code{oracularize_sampler}.  Its theorem contract is \grade{CITED}.
\gt{sec:orac-def}{ground-truth/gt-09-oracularization.tex}
The typed
answer-reduced sampler is the product of this three-role family with the
18-type PCP family.  \code{typed_sampler_product} produces 54 types and
combines levels by a maximum.  The code constructs the product and checks its
shape, but C4b remains \textsc{CONJECTURE}: the latest TB2 critic held promotion
because the lazy \code{PCPCLMap}'s stage depth was not itself enough to make
the full CL decomposition a constructor-level fact.  Local certificate labels
must not be confused with a promoted claim.

The typed decider implements the five guarded checks of
\code{fig:decider-pcp}: global consistency, input consistency, input low degree,
proof encoding (individual and simultaneous low degree), and the PCP game
check.\gt{fig:decider-pcp}{ground-truth/gt-10-answer-reduction.tex}
The Julia names are \code{TypedAnswerReducedDecider},
\code{typed_answer_reduced_decider}, \code{LDParams}, \code{ld_decider}, and
\code{PCPGameCall}.  \code{answer_reduce_pcp} constructs the typed verifier of
level \(\max(\ell,3)\) on its fixture and attaches explicit assumptions.

The full theorem contract is conditional:
\[
\begin{array}{ll}
\textsc{Assume:}&V\text{ is an }\ell\text{-level normal-form verifier},\\
&T(n)=(2^{\lambda n})^\mu,\quad Q(n)=(\lambda n)^\mu,\\
&\Time_D(n)\leq T(n),\quad \Time_S(n)\leq Q(n);\\[2pt]
\textsc{Prove:}&\AnswerReduce(V)\text{ has level }\max\{\ell+2,5\},\\
&\text{the stated polynomial runtimes and directed completeness,}\\
&\text{soundness, and entanglement implications.}
\end{array}
\]
This is \code{thm:ar} in
\code{ground-truth/gt-10-answer-reduction.tex}.  The finite guarded decider is
executable; the quantum implication, detyping, and asymptotic theorem are
\grade{CITED}.

\subsection{The structural hypothesis}

\emph{Introspection} is the cited verifier transformation that makes encoded
prover-side measurement information available to a smaller verifier while
preserving a directed soundness contract.
\gt{thm:introspection}{ground-truth/gt-08-introspection.tex}
\emph{Anchoring} inserts special anchor questions with positive probability;
\emph{parallel repetition} runs several copies with an AND acceptance
condition against possibly correlated strategies.  Their combined repetition
contract is cited from \code{thm:repetition} in
\code{ground-truth/gt-11-parallel-repetition.tex}.

The relevant closure equations are
\[
 \Introspect(\mathcal Q)\subseteq\mathcal Q,
 \qquad \mathsf{PCPCompose}(\mathcal Q)\subseteq\mathcal Q,
 \qquad \mathcal Q\times\mathcal Q\subseteq\mathcal Q.
\]
Concatenation, direct sum, the shared-seed distribution, and product are visibly
CL-preserving in the base inductive datatype.  Affine restriction is linear
after earlier stages choose a direction, and random-point identity tests are
pushforwards of a shared uniform seed.  Typed graph choice is external, and the
cited detyping transformation moves it back into two further CL levels.

This does not characterize every admissible query distribution, prove that a
generic PCPP can be expressed in the algebra, or prove closure under the full
introspection compiler.  The structural statement is exactly
\claim{C7}{CONJECTURE}.  Introspection soundness and normal-form closure are
not implemented.  The present datatype makes several local closure laws
legible, not the global compression theorem.

\section{Correspondence tables}

The evidence grade describes the local certificate leaf.  The claim column is
the current status in \code{claims/CLAIMS.md}; it can remain weaker than a
certificate label when an adversarial verdict has not promoted the claim.
\grade{CONSTRUCTED} means enforced by a constructor,
\grade{CHECKED} means deterministic replay against an attached object, and
\grade{CITED} means imported from the named ground-truth statement.

\subsection{Description front end (planned TB3)}

\setlength{\tabcolsep}{1.5pt}
\scriptsize
\begin{longtable}{P{0.21\textwidth}P{0.16\textwidth}P{0.19\textwidth}P{0.19\textwidth}P{0.12\textwidth}P{0.08\textwidth}}
\toprule
Paper object & Analytic statement & IR sort & Julia name & Grade & Claim \\
\midrule
\endhead
\code{sec:tms}\newline\code{gt-03-prelim.tex}
& Machine conventions in Definition~\ref{def:paper-tm}; costed bridge in
Theorems~\ref{thm:tm-to-l}--\ref{thm:l-to-tm}
& \code{Quoted{A}} & \code{Quote}, \code{Eval}, \code{Specialize}
& \grade{CITED}; absent & none \\
\code{def:decider}\newline\code{gt-05-games-normalform.tex}
& Five-input format in Definition~\ref{def:paper-decider}; bounds preserved by
Theorem~\ref{thm:lambda-preservation} & \code{Quoted{Decider}}
& \code{description_size} & \grade{CITED}; absent & none \\
\code{DESIGN} \S1.1
& Determinism, fuel, quotation, and specialization:
Theorems~\ref{thm:l-determinism}--\ref{thm:quote-eval} and
Lemma~\ref{lem:specialization}
& \code{ClosedProgram}, \code{Quoted} & \code{Eval}, \code{specialize}
& analytic only; absent & none \\
\code{prop:standard-succinct-sat}\newline\code{gt-10-answer-reduction.tex}
& Transition locality (Lemma~\ref{lem:transition-window}) and succinct 3SAT
(Theorem~\ref{thm:l-succinct-sat}) & \code{BoundedTrace}, \code{Succinct3SAT}
& \code{bounded_trace}, \code{cook_levin} & \grade{CITED}; absent & none \\
\code{sec:succinct-deciders}\newline\code{gt-10-answer-reduction.tex}
& Three shared reads become five independent blocks &
\code{SuccinctDecoupled5SAT} & \code{decouple5} & \grade{CITED}; absent & none \\
\bottomrule
\end{longtable}
\normalsize

\subsection{Polynomial rung (TB0)}

\scriptsize
\begin{longtable}{P{0.21\textwidth}P{0.16\textwidth}P{0.19\textwidth}P{0.19\textwidth}P{0.12\textwidth}P{0.08\textwidth}}
\toprule
Paper object & Analytic statement & IR sort & Julia name & Grade & Claim \\
\midrule
\endhead
\code{def:tseitin}\newline\code{gt-10-answer-reduction.tex}
& Circuit becomes a linear-size formula with auxiliary wires &
\code{Circuit}, \code{TseitinFormula} & \code{tseitin}, \code{occurrences}
& \grade{CHECKED} occurrence & C2 \textsc{Conj.} \\
\code{def:formula-arithmetization}\newline\code{gt-10-answer-reduction.tex}
& Boolean connectives become a formal polynomial & \code{Poly}
& \code{arith_q}, \code{actual_degrees} & \grade{CHECKED} & C8 \textsc{Tested} \\
\code{sec:ld-encoding}\newline\code{gt-03-prelim.tex}
& Assignment becomes its multilinear extension & \code{Poly}
& \code{g_a}, \code{dependency_blocks} & \constructed/\newline\grade{Checked}
& C3 \textsc{Tested} \\
\code{prop:zero-basis}\newline\code{gt-10-answer-reduction.tex}
& Divide by \(z_i(1-z_i)\) and require zero remainder &
\code{ZeroDecomposition} & \code{zero_basis_decompose},
\code{verify_zero_decomposition} & \grade{CHECKED} & C2 \textsc{Conj.} \\
\code{fig:pcpverifier}\newline\code{gt-10-answer-reduction.tex}
& Check formula and zero identities at one view & \code{PCPView}
& \code{pcpverifier} & \grade{CHECKED} finite views & C1 \textsc{Conj.} \\
\code{thm:pcp-decider}\newline\code{gt-10-answer-reduction.tex}
& Acceptance \(>1/2\) decodes an accepting trace, assuming low degree &
\code{PCPProof} & \code{build_pcp}, \code{parameter_policy}
& \grade{CITED}/\newline\grade{Checked} parts & C5 \textsc{Sketch} \\
\bottomrule
\end{longtable}
\normalsize

\subsection{CL and low-degree-test rung (TB1)}

\scriptsize
\begin{longtable}{P{0.21\textwidth}P{0.16\textwidth}P{0.19\textwidth}P{0.19\textwidth}P{0.12\textwidth}P{0.08\textwidth}}
\toprule
Paper object & Analytic statement & IR sort & Julia name & Grade & Claim \\
\midrule
\endhead
\code{def:cl-func}\newline\code{gt-04-cl.tex}
& CL functions are inductive conditional linear stages & \code{AbstractCL}
& \code{CLZero}, \code{CLStep}, \code{level} & \constructed
& C4a \textsc{Tested} \\
\code{lem:cl-kth}\newline\code{gt-04-cl.tex}
& A marginal is the sum of the first \(k\) stages & \code{CLMarginal}
& \code{marginal_k}, \code{sum_stage_outputs} & \grade{CHECKED}
& C4a \textsc{Tested} \\
\code{lem:cl-concat}, \code{lem:cl-func-prod}\newline\code{gt-04-cl.tex}
& Concatenation adds levels; direct sum takes a maximum & \code{AbstractCL}
& \code{concatenate}, \code{direct_sum}, \code{product}
& \constructed & C7 \textsc{Conj.} \\
\code{lem:alnf}, \code{lem:dlnf}\newline\code{gt-07-ldt.tex}
& CL pushforwards equal the line/point distributions & \code{CLDistribution}
& \code{L_Point}, \code{L_ALine}, \code{L_DLine}, \code{histogram}
& \grade{CHECKED} TB1 & C4a \textsc{Tested} \\
\code{fig:ld-decider}\newline\code{gt-07-ldt.tex}
& Point/line answer consistency and degree format & \code{LDParams}
& \code{ld_decider}, \code{restrict} & \grade{CHECKED} fixture & C4a \textsc{Tested} \\
\code{lem:ld-soundness}\newline\code{gt-07-ldt.tex}
& Successful quantum strategy yields polynomial measurements & theorem leaf
& no proof checker & \grade{CITED} & C5 \textsc{Sketch} \\
\bottomrule
\end{longtable}
\normalsize

\subsection{Typed answer-reduction rung (TB2)}

\scriptsize
\begin{longtable}{P{0.21\textwidth}P{0.16\textwidth}P{0.19\textwidth}P{0.19\textwidth}P{0.12\textwidth}P{0.08\textwidth}}
\toprule
Paper object & Analytic statement & IR sort & Julia name & Grade & Claim \\
\midrule
\endhead
\code{def:typed-sampler}\newline\code{gt-06-types.tex}
& A type graph selects a pair of CL maps sharing one seed & \code{TypedSampler}
& \code{TypedSampler}, \code{sample}, \code{pad_level}
& \constructed\ base & C4b \textsc{Conj.} \\
\code{sec:ld-compiler}\newline\code{gt-10-answer-reduction.tex}
& Six low-degree copies form an 18-type family & \code{PCPCLMap}
& \code{pcp_sampler}, \code{intrinsic_pcp_levels}
& local \grade{CHECKED}; held & C4b \textsc{Conj.} \\
\code{sec:ld-compiler}\newline\code{gt-10-answer-reduction.tex}
& Product with the oracularized sampler has 54 types &
\code{TypedProductCL} & \code{oracularize_sampler},
\code{typed_sampler_product} & \grade{CHECKED} shape; theorem cited
& C4b \textsc{Conj.} \\
\code{fig:decider-pcp}\newline\code{gt-10-answer-reduction.tex}
& Five guarded checks connect input, proof, low degree, and PCP &
\code{TypedAnswer}\newline\code{ReducedDecider} & \code{typed_answer_reduced_decider}
& \grade{CHECKED} finite paths & no current row \\
\code{lem:detyping-verifiers}\newline\code{gt-06-types.tex}
& Remove types at cost \(+2\) levels and \(16^{54}\) soundness factor &
\code{CitedDetyped}\newline\code{Verifier} & \code{detype} & \grade{CITED} & C4b \textsc{Conj.} \\
\code{thm:ar}\newline\code{gt-10-answer-reduction.tex}
& Conditional complexity, completeness, soundness, entanglement contract &
\code{TypedAnswer}\newline\code{ReducedVerifier} & \code{answer_reduce_pcp},
\code{AnswerReduce} & \grade{Checked}/\newline\grade{Cited} & C5 \textsc{Sketch} \\
\bottomrule
\end{longtable}
\normalsize

\subsection{Midpoint diagnostic (TB0.5)}

\scriptsize
\begin{longtable}{P{0.21\textwidth}P{0.16\textwidth}P{0.19\textwidth}P{0.19\textwidth}P{0.12\textwidth}P{0.08\textwidth}}
\toprule
Analytic anchor & Analytic statement & IR sort & Implementation & Grade & Claim \\
\midrule
\endhead
\code{handoff.md}\newline midpoint diagnostic
& False level-\(n\) claim has value \(1-2^{-n}\) under the stated orbit condition
& \code{Test/Ask/Coin} terms & \code{toys/midpoint}
& proof plus exact tests & C6 \textsc{Proved} \\
\code{toys/midpoint/PROOF.md}
& Sequential \(r\)-copy AND value is \((1-2^{-n})^r\); constant gap costs
\(\Theta(2^n)\) copies & exact dynamic program & \code{sequential_value}
& proof plus exact tests & N1 \textsc{Proved} \\
\bottomrule
\end{longtable}
\normalsize

\subsection{Compression and fixed point (planned TB4)}

\scriptsize
\begin{longtable}{P{0.21\textwidth}P{0.16\textwidth}P{0.19\textwidth}P{0.19\textwidth}P{0.12\textwidth}P{0.08\textwidth}}
\toprule
Paper object & Analytic statement & IR sort & Julia name & Grade & Claim \\
\midrule
\endhead
\code{fig:compress}\newline\code{gt-12-compression.tex}
& \(\Repeat\circ\AnswerReduce\circ\Introspect\) &
\code{ClosedProgram}\newline\code{{Compressor}} & \code{Compress} & \grade{CITED}; absent
& C7 \textsc{Conj.} \\
\code{thm:compression}\newline\code{gt-12-compression.tex}
& Nine-level gap-preserving compression with dimension bound & contract tree
& no current implementation & \grade{CITED} & C7 \textsc{Conj.} \\
\code{fig:halt_f}\newline\code{gt-12-compression.tex}
& Self-application, Kleene, and YCode coincide by
Theorems~\ref{thm:kleene-two}--\ref{thm:fixed-point-equivalence} &
\code{PartialProgram}, \code{Quoted} & \code{Hole}, \code{Specialize},
\code{Fix}/\code{YCode} & \grade{CITED}; absent & none \\
\bottomrule
\end{longtable}
\normalsize

\section{What is analytic and what is computed}

\subsection{The evidence boundary}

The analytic content is the sequence of transformations and the implications
attached to them.  The computer currently realizes a strict middle segment.

\begin{center}
\small
\begin{tabularx}{\textwidth}{P{0.20\textwidth}P{0.35\textwidth}X}
\toprule
Class & Present content & What would be required to move the boundary \\
\midrule
Executed constructions
& Circuit/Tseitin syntax; formula arithmetization; multilinear extensions;
\(c_0\); zero-basis quotients; point/line maps; typed products; guarded TB2
decider
& Scale-independent front end for quoted programs and a constructor-faithful
lazy CL stage representation for the PCP sampler \\
Certified identities
& Normalized coefficient equality; zero remainder; structural and actual
degrees; dependency blocks; CL marginals; TB1 exact histograms; finite parser
and branch checks
& General semantic certificates connecting trace, SAT, Tseitin, and
arithmetization, plus promotion through adversarial verdicts \\
Cited analytic leaves
& Cook--Levin asymptotics; general Tseitin semantics; PCP soundness;
low-degree extraction; oracularization; detyping; introspection; rigidity;
anchored repetition; compression
& Formal proofs over quantified machines, distributions, Hilbert spaces,
measurements, approximation metrics, and asymptotic bounds---not more finite
samples \\
\bottomrule
\end{tabularx}
\end{center}

Lean or another proof assistant could check the algebraic core once finite
fields, sparse polynomials, interpolation, and polynomial division are
formalized.  That would move the general zero-basis identity and degree lemmas
from prose to proof terms.  It would not by itself certify
\code{lem:ld-soundness}: that leaf requires a library for finite-dimensional
Hilbert spaces, POVMs, state-dependent distances, tensor-code tests, and the
quantitative extraction argument.  Rigidity and oracularization require the
same operator layer.  Detyping additionally needs a formal probability and
rejection-sampling argument with the \(16^{|\mathrm{Type}|}\) loss.

Anchoring and parallel repetition are combined in \(\Repeat\) and are
theorem-backed rather than executed here.  Introspection and compression would
require a still broader object: a
formal semantics for quoted verifier descriptions, resource bounds under
specialization and universal evaluation, and a proof that the directed value
and entanglement implications compose.  The relevant final statements are
\code{fig:compress} and \code{thm:compression} in
\code{ground-truth/gt-12-compression.tex}.

\subsection{The midpoint diagnostic}

The recursive midpoint verifier for \(y=f^{2^n}(x)\) asks for \(z\) and checks
one of
\[
 z=f^{2^{n-1}}(x),\qquad y=f^{2^{n-1}}(z)
\]
with a fresh fair coin.  It has a compact fixed-point description and perfect
completeness.  Under the orbit-prefix condition stated in C6, the exact optimum
for a false claim is nevertheless
\[
 p_n=1-2^{-n}.
\]
Claim C6 is \textsc{PROVED}.  For adaptive sequential AND repetition, N1 is
also \textsc{PROVED} and gives value \((1-2^{-n})^r\); reaching error at most
\(1/2\) requires \(r=\Theta(2^n)\), hence \(\Theta(n2^n)\) transcript cost in
the stated unit-cost model.  Neither claim concerns quantum parallel repetition
or the actual \(\Compress\) operator.

The diagnostic isolates the point of the whole construction.  A recursive
term can represent an exponentially long computation, and a fixed point can
make that term self-referential, while the acceptance gap collapses too quickly
for iteration.  Representability is easy; an iteratable, gap-preserving
analytic transformation is the hard requirement.

\subsection{Final accounting}

The lambda reformulation removes none of the mathematically difficult leaves.
It makes quotation, hardwiring, description size, timeout, and self-reference
explicit in one syntax.  The polynomial IR makes the Cook--Levin-to-PCP middle
visible as actual algebra, and the CL datatype exposes several sampler closure
laws.  What remains unchanged is the proof that high success in the typed
quantum game yields consistent low-degree polynomial measurements, survives
oracularization and detyping with the stated quantitative loss, survives
introspection and anchored parallel repetition, and composes into the
gap-preserving compression theorem.  Those are analytic theorems, presently
\grade{CITED}; they are neither consequences of homoiconicity nor conclusions
of the finite computations reported here.

\end{document}
